\documentclass[11pt,letterpaper]{article}
\usepackage[margin=1in]{geometry}
\usepackage[T1]{fontenc}
\usepackage{lmodern}
\usepackage{amsmath,amssymb,amsthm,mathtools}
\usepackage{booktabs,array,tabularx}
\usepackage{microtype}
\usepackage[numbers,sort&compress]{natbib}
\usepackage{xcolor}
\usepackage[colorlinks=true,linkcolor=blue!45!black,citecolor=blue!45!black,urlcolor=blue!45!black]{hyperref}
\hypersetup{pdftitle={Positional Information for Budget-Adaptive Markov Prediction},pdfauthor={},pdfsubject={Budget-adaptive Markov prediction and finite positional statistics}}
\newcommand{\TV}{\operatorname{TV}}
\newcommand{\E}{\mathbb E}
\newcommand{\ind}{\mathbf 1}
\newcommand{\Prb}{\mathbb P}
\newcommand{\eps}{\varepsilon}
\newcommand{\cT}{\mathcal T}
\newcommand{\cK}{\mathcal K}
\newcommand{\R}{\mathbb R}
\DeclareMathOperator{\diag}{diag}
\DeclareMathOperator{\tr}{tr}
\DeclareMathOperator{\Var}{Var}
\DeclareMathOperator{\Cov}{Cov}
\DeclareMathOperator{\KL}{KL}
\newtheorem{theorem}{Theorem}[section]
\newtheorem{proposition}[theorem]{Proposition}
\newtheorem{lemma}[theorem]{Lemma}
\newtheorem{corollary}[theorem]{Corollary}
\theoremstyle{definition}
\newtheorem{definition}[theorem]{Definition}
\theoremstyle{remark}
\newtheorem{remark}[theorem]{Remark}
\numberwithin{equation}{section}
\setlength{\emergencystretch}{2em}
\title{Positional Information for Budget-Adaptive Markov Prediction}
\author{}
\date{3 November 2026}

\begin{document}
\maketitle
\begin{abstract}
We study next-symbol prediction from one binary order-two Markov trajectory after an unknown number of symbol edits, with the terminal pair protected. Loss is expected total variation at that actual pair. At a supplied all-state two-step minorization certificate $1/M$, the known-budget risk scale is $g_M(k)=\max\{M^{-1/2},k/M\}$. For every $M\ge64$ and $2\le r\le M+2$, we prove that the optimal competitive factor of one budget-blind map observing exact triple counts and the first and last $r$ symbols has order $1+\log((M+2)/r)$, with explicit universal constants. The upper uses only counts and the suffix; the lower grants both windows and compares the same map across budgets. A whole-word scan can instead supply one exact positional gate bit beyond counts and the protected pair. We then give a perturbation bound for approximate full and half-prefix counts, a fully accounted finite masked-attention/Boolean-ReLU realization, and a matching clean precision bound for its twelve prescribed rounded coordinates. These latter results calibrate a specific observation interface: they do not establish neural superiority over same-information counters, a learning result, or a precision law for an entire network. Complete proofs and the distinctions among the observation models are included.
\end{abstract}

\section{The observation question}
A transition histogram records which local patterns occurred but not when they occurred. For a clean homogeneous Markov trajectory, this is consistent with the usual likelihood factorization into transition counts and the initial state. With position-sensitive corruption, however, the observation channel changes. A predictor may need to distinguish an empirical row supported by older visits from one populated almost entirely near the prediction time. The question here is how much positional information is needed when the number of edits is not supplied.

We study one concrete experiment: a binary order-two source, $M$ observed departures, and an unmodified terminal pair. The source obeys an all-state two-step minorization with supplied coefficient $1/M$. An offline kernel may change at most $k$ earlier symbols on every realization. The predictor sees exact triple counts and prescribed boundary windows, but not $k$. Its risk is expected unsquared total variation to the next-symbol distribution at the actual terminal pair. In this experiment, giving $k$ to the predictor makes counts and the protected pair sufficient for the optimal \emph{rate}; no claim of statistical sufficiency is needed.

Our main result characterizes the competitive factor simultaneously for every finite horizon and boundary width: its order is $1+\log((M+2)/r)$ with universal constants. One map at each fixed horizon must serve every budget. The positive lower limiting fraction criterion and logarithmic fixed-fraction price follow at the order level; a separate fixed-fraction proof supplies a stronger explicit lower constant. A separate corollary shows why raw access and computed messages must not be identified: a whole-word scan can produce a single exact gate bit that suffices beyond counts. The scan, exact counters, and protected query are part of that corollary's information model.

The finite implementation is subordinate to this information question. A discontinuous accept/reject gate need not be recovered exactly: uniform approximation of its full and half-prefix statistics controls expected prediction error. An explicit masked-attention program realizes these statistics, while a classical one-pass counter realizes the same information more simply. We keep these computational comparisons separate from the restricted-observation lower bound.

The proof mechanisms are complementary. For the upper bound, a gate retains evidence of visits before a suffix, removing the shortest-return terms from a harmonic convolution. For the lower bound, clean alternatives with small prediction risk are transformed into legally edited null words whose restricted observations match those alternatives. The clean guarantee is used at the same horizon and supplied certificate as the attacked guarantee. Complete arguments appear in the appendices. Sections~\ref{sec:experiment}--\ref{sec:bit} give the statistical comparison, Section~\ref{sec:finite-implementation} calibrates one finite implementation, and Section~\ref{sec:related} states the relationship to prior work.

\section{Experiment and benchmark}\label{sec:experiment}
All logarithms are natural. Let $M\ge64$ be an integer and
\[
 X=(X_0,\ldots,X_{M+1}),\qquad Y_i=(X_i,X_{i+1}),\quad 0\le i\le M.
\]
Thus there are $M+2$ raw symbols, $M$ departures, and the terminal pair $Y_M$. A homogeneous order-two transition table $\theta_{ab}(c)$ induces the four-state pair kernel
$P((a,b),(b,c))=\theta_{ab}(c)$. There is a qualifying stationary law $\Pi$ such that
\begin{equation}\label{eq:certificate}
 \Pi P=\Pi,\qquad P^2(u,v)\ge\beta\Pi(v)\quad\text{for all }u,v\in\{00,01,10,11\}.
\end{equation}
The condition includes rows outside stationary support. Throughout the main results the supplied certificate is \emph{exactly} $\beta=1/M$; it need not be the source's largest possible coefficient. Neither $\theta$ nor $\Pi$ is revealed. Stationary initialization means $Y_0\sim\Pi$; arbitrary initialization means any hidden law $\mu$, including laws supported outside the stationary support.

\paragraph{Strict edits and actual-terminal loss.}
For $0\le k\le M$, a legal randomized kernel $K_{M,k}$ maps full binary words to full binary words. On every input $x$, every output $z$ in its support satisfies
\begin{equation}\label{eq:attack}
 z_M=x_M,\qquad z_{M+1}=x_{M+1},\qquad
 \sum_{i=0}^{M-1}\ind\{z_i\ne x_i\}\le k.
\end{equation}
The kernel is one fixed map on words, not supplied the true source table, stationary or initial law, or predictor's random seed. It may inspect the whole input and use public parameters and $k$; lower-bound kernels may hardcode a predeclared auxiliary family. No causal or computational restriction is imposed. Its randomness is independent of the predictor's. This formulation specifies the common kernels used in the proofs; source independence alone is not asserted to distinguish all standard contamination models.

For $Z\sim K_{M,k}(X,\cdot)$ and an estimator $A$ based on observation $\cT(Z)$, define
\begin{equation}\label{eq:risk}
 R(A;\theta,\mu,K_{M,k})
 =\E\,\TV\bigl(A(\cT(Z)),\theta_{Y_M}\bigr).
\end{equation}
The expectation includes the path and both independent randomizations. The target is the source's next-symbol law at the \emph{actual protected pair}, not an independent query. For binary distributions, TV is absolute error in either coordinate. Risk is unconditional.

\paragraph{Observed information.}
For $a,b,c\in\{0,1\}$ let
\[
 H_{abc}(z)=\sum_{i=0}^{M-1}\ind\{(z_i,z_{i+1},z_{i+2})=(a,b,c)\}.
\]
For prescribed $2\le r\le M+2$, the observation is
\begin{equation}\label{eq:observation}
 \cT_r(z)=\bigl((H_{abc}(z))_{a,b,c},\ z_0,\ldots,z_{r-1},\
 z_{M+2-r},\ldots,z_{M+1}\bigr).
\end{equation}
The windows may overlap. Counts and raw symbols are exact; neither edit locations nor the clean interior are supplied. In row $u=(a,b)$ use the designated successor $d(u)=1-a$, and define
\[
 J'_u=H_{ab0}(Z)+H_{ab1}(Z),\qquad C'_u=H_{ab,d(u)}(Z),\qquad
 q_u=\theta_{ab}(d(u)).
\]
Primes indicate corrupted counts. Unprimed counts are clean analysis variables. At the selected row $u=Y_M$, subscripts are omitted. A ratio $C/J$ is evaluated only when $J>0$.

\paragraph{One family for all budgets.}
Use the explicitly normalized rate profile
\begin{equation}\label{eq:benchmark}
 g_M(k)=\max\{M^{-1/2},k/M\},\qquad 0\le k\le M.
\end{equation}
For a prescribed sequence $r_\bullet=(r_M)_{M\ge64}$, let $C_{\rm stat}(r_\bullet)$ be the infimum of constants $C\ge0$ for which one estimator family $(A_M)$, based on $\cT_{r_M}$, satisfies
\begin{equation}\label{eq:factor}
 R(A_M;\theta,\Pi,K_{M,k})\le Cg_M(k)
\end{equation}
for every $M\ge64$, every $0\le k\le M$, every qualifying source and every legal kernel. Define $C_{\rm all}$ by additionally quantifying over every $\mu$. The family may depend on the prescribed sequence, $M$, $r_M$, and supplied $1/M$, but not on $k$ or hidden source quantities. The infimum need not be attained; an empty admissible set has value $+\infty$.

\paragraph{Finite-horizon competitive factors.}
Separately from the sequence factors $C$, fix $(M,r)$ and a single
$k$-blind map $A$ based on $\cT_r$. Define
\begin{equation}\label{eq:finite-factor}
 \Gamma_{\rm stat}(A;M,r)
 =\sup_{\substack{0\le k\le M,\ \theta,\Pi,K_{M,k}\\
                  \text{qualifying stationary experiments}}}
 \frac{R(A;\theta,\Pi,K_{M,k})}{g_M(k)},\qquad
 \Gamma_{\rm stat}(M,r)=\inf_A\Gamma_{\rm stat}(A;M,r).
\end{equation}
Define $\Gamma_{\rm all}(A;M,r)$ and $\Gamma_{\rm all}(M,r)$ by also
allowing every initial law $\mu$. The map may depend on $M,r$ and the
supplied $1/M$, but the very same map serves all $k$ at that horizon.
These are finite-horizon infima, not the sequence quantities $C$ above;
no infimum is assumed attained. Every map has factor at most $\sqrt M$
because TV is at most one.

The comparison is to a genuine same-information known-budget scale, rather than to an arbitrary denominator. Write $R^{\cT}_{\rm stat}(M,k)$ and $R^{\cT}_{\rm all}(M,k)$ for the respective minimax risks when $k$ is also supplied.
\begin{proposition}[Known-budget scale]\label{prop:oracle}
For the full corrupted history, for $\cT_r$, and for counts plus the protected pair,
\begin{equation}\label{eq:oracle}
 \frac{g_M(k)}{96}\le R^{\cT}_{\rm stat}(M,k)
 \le R^{\cT}_{\rm all}(M,k)\le191488g_M(k).
\end{equation}
\end{proposition}
The upper is count-only and allows arbitrary starts; the lower holds even for full histories. Thus a factor $C$ relative to $g$ gives a factor at most $96C$ relative to the corresponding minimax oracle. Conversely, a factor $D$ relative to that oracle gives a factor at most $191488D$ relative to $g$. Proofs are in Lemma~\ref{lem:oracle-floor} and Proposition~\ref{prop:count-oracle}.

For comparison with a more general certificate-dependent profile, define $F_2(M,0,\beta)=0$ and, for $k\ge1$,
\[
 F_2(M,k,\beta)=\frac{k}{M}\left[1+\log\left(1+
 \frac1{k\max\{\sqrt\beta,1/M\}}\right)\right],\qquad
 H_2=\min\{1,\max\{M^{-1/2},F_2\}\}.
\]
At the certificate used here, $\log(1+x)\le x$ gives $g_M(k)\le H_2(M,k,1/M)\le2g_M(k)$. Hence the two normalizations change optimal factors by at most two; all constants below are stated directly for $g$.

\section{Finite-horizon characterization and its consequences}\label{sec:results}
\begin{theorem}[Simultaneous finite-horizon window price]\label{thm:finite-window}
For every integer $M\ge64$ and $2\le r\le M+2$, put
$\ell=1+\log((M+2)/r)$. Then
\begin{equation}\label{eq:finite-window}
 2^{-30}\ell\le\Gamma_{\rm stat}(M,r)
 \le\Gamma_{\rm all}(M,r)\le2^{20}\ell.
\end{equation}
The upper uses only exact counts and the suffix. The lower grants both
boundary windows. All experiments use the supplied certificate $1/M$.
\end{theorem}
The proof in Section~\ref{sec:finite-window-proof} checks the complete finite
domain of the matched-observation bound. It tests the clean risk of the same
map before selecting an attacked budget; it is not a logarithmic risk lower
at each budget. Linear widths have constant factor; for fixed $0<\gamma<1$,
$r=\max\{2,\lceil M^\gamma\rceil\}$ has order $\log M$, as does $r=2$.
The constants are conservative, not a claim of sharp practical thresholds.

The next sequence statements are consequences at the order level. Their
original separate proofs are retained in
Proposition~\ref{prop:lower-consequences}, including the stronger
fixed-fraction lower constant.

\begin{theorem}[Exact raw-window criterion]\label{thm:width}
For every prescribed sequence of integers $2\le r_M\le M+2$,
\[
 C_{\rm stat}(r_\bullet)<\infty
 \quad\Longleftrightarrow\quad C_{\rm all}(r_\bullet)<\infty
 \quad\Longleftrightarrow\quad
 \liminf_{M\to\infty}\frac{r_M}{M}>0.
\]
No monotonicity of $r_M$ is required. The factor may depend on the prescribed sequence.
\end{theorem}
Necessity already uses stationary witnesses. Sufficiency includes unsupported starts. A finite initial set of horizons is harmless because $g_M(k)\ge M^{-1/2}$. The conclusion is about one budget-blind family, not a separate minimax estimator chosen for each $k$.

\begin{theorem}[Price of a fixed fraction]\label{thm:fraction}
Fix $0<\alpha\le1/4096$ and $r_M=\max\{2,\lceil\alpha M\rceil\}$. Abbreviating the corresponding factors by $C_{\rm stat}(\alpha)$ and $C_{\rm all}(\alpha)$,
\begin{equation}\label{eq:fraction}
 2^{-26}\log\frac1{32\alpha}
 \le C_{\rm stat}(\alpha)\le C_{\rm all}(\alpha)
 \le2^{19}\left(1+\log\frac2\alpha\right).
\end{equation}
The upper requires only the histogram and suffix; the lower permits both boundary windows.
\end{theorem}
Thus the optimal uniform factor has order $\log(1/\alpha)$ as $\alpha\downarrow0$. The upper is uniform in every $M,k$ for each fixed $\alpha$. The lower first sends $M\to\infty$ with $\alpha$ fixed and with a proposed finite factor fixed. This fixed-fraction proof takes limits in the stated order. The finite characterization above requires no limit in either parameter; neither result is a logarithmic risk lower bound at every budget. The constants are explicit but not optimized. Proposition~\ref{prop:lower-consequences} proves Theorems~\ref{thm:width} and~\ref{thm:fraction} from a finite upper bound and a finite matched-observation lower bound.

\subsection{An explicit estimator}
Set
\begin{equation}\label{eq:suffix-length}
 m=\min\{r-2,\lfloor M/2\rfloor\},\qquad N=M-m,\qquad
 L=1+\log\frac{M}{m+1}.
\end{equation}
The last $m+2$ symbols reveal exactly the $m$ departures indexed $M-m,\ldots,M-1$ and the terminal pair. Let $T'_u$ count visits to $u$ at those departures and put $O'_u=J'_u-T'_u$. Then $O'_u$ is the count on indices $0,\ldots,N-1$; it is a direct old-prefix count even though it is computed by subtraction. At the selected row, predict designated probability
\begin{equation}\label{eq:gate}
 \widehat q=
 \begin{cases}
 C'/J',& J'>0\ \text{and}\ 512O'\ge J',\\
 0,&\text{otherwise}.
 \end{cases}
\end{equation}
The complementary symbol receives probability $1-\widehat q$. This map does not use $k$ or $\beta$. When $r=2$, $m=0$ and $O'=J'$; the terminal pair is not a departure. Using $r-1$ in place of $r-2$ would overcount the available information.

The finite upper bound at the main certificate, proved in Proposition~\ref{prop:finite-upper}, is
\begin{equation}\label{eq:finite-upper}
 R\le\min\left\{1,\frac{131072}{\sqrt M}
 +(8448L+147744)\frac{k}{M}\right\}
 \le287264L g_M(k).
\end{equation}
If $r_M/M$ is eventually bounded below, then $L$ is eventually bounded. For fixed $\alpha$ as in Theorem~\ref{thm:fraction}, $m+1\ge\alpha M/2$, including the horizons at which $r_M=2$. This proves the two stated upper conclusions without a normalization loss.

\section{Why positions enter the proof}\label{sec:mechanisms}
\paragraph{Older visits and a truncated harmonic sum.}
The gate is not an estimate of the budget. It requires a positive fraction of the selected row's visits to precede the observed suffix, not merely one old occurrence. This proportional balance makes gate-switch errors chargeable to the edit budget; one forged old occurrence alone must not unlock a large corrupted suffix count. For a clean $d$-state homogeneous chain and any start, the relevant occupation inequality is
\begin{equation}\label{eq:reciprocal-main}
 \E\frac{\ind\{O\ge cJ\}}{J+1}
 \le(3d^2+2d^3)\frac{1+\log(M/(m+1))}{M},\qquad c>0.
\end{equation}
Here $J$ is the number of departures from the actual terminal state and $O$ counts those before time $M-m$. Lemma~\ref{lem:old-inverse} proves this by an exact first-hit convolution and a reciprocal-return identity. In this analytic bound the event can be enlarged to the existence of an old visit. That enlargement removes short backward-time terms; the surviving harmonic range produces the logarithm. The implemented gate still requires proportional balance. Generic finite-state first-hit pointmass estimates have precedents in \citet{NorrisPeresZhai2014}; the required gated convolution is proved explicitly.

A strict $k$-edit outcome changes each relevant count by at most $6k$. A deterministic true-gate/false-gate split therefore reduces the error to a clean empirical comparator, the reciprocal count in \eqref{eq:reciprocal-main}, a weighted imbalance term, and two small-occupation events. The all-state certificate controls the imbalance even from unsupported starts. No conditional independence of departures given the random terminal row or its occupation is assumed. The supplement's row-stack argument instead starts from independent row-wise randomness and then bounds the selected empirical deviation.

\paragraph{Mixture domination and legal observation matching.}
The lower bound uses stationary alternatives whose designated exit probabilities are
\[
 (q_{00},q_{01},q_{10},q_{11})=(a,1/2,1/2,t),\qquad
 \Pi_a=\frac{(t/a,2t,2t,1)}{1+t/a+4t}.
\]
An integral over $a$ creates the logarithmic term. Conditioning on a fixed departure count, one counts marked excursions and holding slots; the resulting occupation-mixture bound controls each proposed output mass. An interior loop relocation preserves the alternative's triple histogram and both of its observed windows, while making it close to a simple two-run null word. These are the alternative's windows, not windows of the unedited null input: changing an observed input-window symbol is legal but must be charged to the budget.

The attack normalizes only nonzero proposal masses, samples once, performs the relocation, and checks the actual Hamming cost and terminal pair. If any check fails, it returns its input unchanged. It is therefore a common source-independent kernel satisfying the budget on every outcome, not only on a high-probability event. In the risk comparison, nonnegative loss is discarded before proposal-mass domination; a signed triangle inequality is used only afterwards. Lemmas~\ref{lem:mixture} and~\ref{lem:relocation} and Proposition~\ref{prop:finite-lower} supply these steps.

For sublinear widths along any subsequence, a budget between the window and horizon scales makes the clean risk negligible relative to the alternative's exit scale while the mixture logarithm diverges. For fixed $\alpha$, a budget proportional to the window yields $\log(1/\alpha)$. The alternatives may depend on this tested budget, but their clean risk is bounded using the same estimator at the same $M$ and supplied $1/M$. This constrained-risk logic is related to general risk-transfer arguments such as \citet{DuchiRuan2021}; the dependent-word construction and legal observation matching are additional obligations.

\section{One computed positional bit}\label{sec:bit}
Now change the access model. The decoder receives exact triple counts and the protected pair, optionally also the initial pair. An encoder may scan the entire already corrupted word and send one symbol in a fixed finite alphabet $\mathcal A$. It knows $M$ but not $k$, $\beta$, or the source. The message is exact and fault-free after construction. There are no timing, framing, variable-length, or other side channels, and no shared source-dependent information.

\begin{corollary}[Minimum supplemental alphabet]\label{cor:onebit}
In this scan-and-message model, the minimum $|\mathcal A|$ for a finite uniform factor relative to $g_M(k)$ at supplied $1/M$ is two. This holds for stationary or arbitrary initialization, with or without the initial pair supplied.
\end{corollary}
For the upper, put $n=\lfloor M/2\rfloor$, select $u=Y_M$, and send
\begin{equation}\label{eq:bit}
 B=\ind\{512J'_1\ge J'\},\qquad
 J'_1=\sum_{i=0}^{n-1}\ind\{(Z_i,Z_{i+1})=u\}.
\end{equation}
The decoder outputs $C'/J'$ if $J'>0$ and $B=1$, and zero otherwise. Equality is accepted. If $J'=0$, the bit equals one, but the explicit positive-denominator test enforces the fallback. Proposition~\ref{prop:finite-upper} gives all-start risk at most
$\min\{1,2^{17}M^{-1/2}+2^{23}k/M\}$.
A one-symbol alphabet supplies no information. Theorem~\ref{thm:width} with $r_M=2$ rules it out even when both endpoint pairs are present. This proves minimality without another lower-bound construction.

\paragraph{A finite word illustration, not a risk experiment.}
For $M=65$, the two words in Table~\ref{tab:words} have the same initial pair $01$, terminal pair $10$, and triple counts, in lexicographic order,
\[
 (H_{000},H_{001},H_{010},H_{011},H_{100},H_{101},H_{110},H_{111})
 =(0,1,1,1,1,0,1,60).
\]
At row $10$, $J'=C'=1$. The single departure from that row lies in different halves, so the computed gate differs although $\cT_2$ is identical. The example verifies a combinatorial collision only: it gives no source probabilities, statistical risk comparison, or lower-bound proof.
\begin{table}[ht]
\centering
\caption{An exact $67$-symbol collision. The notation $1^{62}$ means $62$ consecutive ones; the first half contains departures $0$ through $31$.}\label{tab:words}
\begin{tabular}{lccc}
\toprule
Word & $J'_1$ & $B$ & Designated prediction\\
\midrule
$0100\,1^{62}\,0$ & $1$ & $1$ & $1$\\
$0\,1^{62}\,0010$ & $0$ & $0$ & $0$\\
\bottomrule
\end{tabular}
\end{table}

This is a worst-case supplemental alphabet result, not an entropy or mutual-information statement. The baseline exact counts are uncharged inputs. One computed bit is not one observed raw coordinate and does not remove the whole-word scan. A counter-based implementation is a natural baseline: maintain the eight triple counts and four first-half row counts, retain the endpoint, and select its row after the scan. This uses $O(\log M)$ counter workspace, $O(M)$ counter updates, and at most $O(M\log M)$ naive bit operations. No neural mechanism is needed for this statistical rate guarantee.

\section{Approximate statistics and a finite realization}\label{sec:finite-implementation}
The positional gain is a property of information, not of an attention architecture. We now fix one finite representation of that information and distinguish three questions: stability of the statistical decoder, implementation of the representation, and the precision required by an arbitrary decoder of that fixed representation. Throughout this section $h$ is bank precision and $b$ is output precision; $r$ retains its meaning as raw-window width.

\subsection{A direct full and half-prefix bank}
For every row $u$, retain the full counts $J'_u,C'_u$ already defined, and let
\[
 D'_u=\sum_{i=0}^{\lfloor M/2\rfloor-1}
       \ind\{(Z_i,Z_{i+1})=u\}.
\]
Thus $D'_u/M$ is a subfrequency normalized by $M$, not by the half length. The twelve direct coordinates are
\begin{equation}\label{eq:direct-bank}
 (J'_u/M,C'_u/M,D'_u/M)_{u\in\{00,01,10,11\}}.
\end{equation}
The decoder additionally receives the exact protected terminal pair. No other historical information is implicit in this bank.

Suppose a computation approximates each unnormalized coordinate within $e\ge0$ on every allowed outcome. The errors may be correlated, real-valued, and inconsistent with any word. Clip the total to $[0,M]$, then the successor and half count to the resulting interval $[0,\bar J_u]$. At the selected row, use
\begin{equation}\label{eq:approx-gate-main}
 \widehat q=\begin{cases}
 \bar C/\bar J,&\bar J>0\text{ and }512\bar D\ge\bar J,\\
 0,&\text{otherwise}.
 \end{cases}
\end{equation}
Sequential clipping preserves coordinate error, including when the total rounds to zero. Theorem~\ref{thm:approximate-gate} proves that any further probability approximation within $\xi$ yields, at the critical certificate,
\begin{equation}\label{eq:approx-risk-main}
 R\le\min\left\{1,\frac{131072}{\sqrt M}
       +6433056\frac{k}{M}+1072176\frac eM+\xi\right\}.
\end{equation}
The single decoder is independent of $k,e$, and $\xi$. The appendix also states the general positive-certificate bound.

This stability statement does not require a margin separating the gate from equality. On clean rows of large occupation, an approximate false acceptance or rejection is incompatible with the gap between the analytic balance thresholds $1/1024$ and $1/256$. The remaining gate changes occur on small clean rows and are paid by the actual-terminal occupation tail. The proof handles real error radii with integer ceilings; it does not condition the row departures into independent samples.

\subsection{A paid masked-attention implementation}
Set $Q=2^h$ with $h\ge0$. Round each coordinate of \eqref{eq:direct-bank} separately to the nearest point of $\{0,1/Q,\ldots,1\}$, with ties upward. Let its integer numerator be $B_i\in\{0,\ldots,Q\}$. The represented endpoint is exact. The allowed numeric fault replaces the twelve master numerators jointly by canonical integers $B_i^*$ in the same range, satisfying
\[
 \max_i |B_i^*/Q-B_i/Q|\le\delta.
\]
Only those replaced values, the endpoint, and public source-independent format metadata are available downstream. Endpoints, masks, addresses, and the subsequent arithmetic remain exact. This is a post-rounding coordinate replacement model, not arbitrary analog noise or register bit flips.

Theorem~\ref{thm:finite-realization} constructs the bank for every binary word. Two zero-score singleton-mask heads copy the preceding one and two bits. At departure key $i+2$, tokenwise Boolean/ReLU circuits form row, designated-successor, and first-half indicators. A third zero-score head at a synthetic final query averages exactly the $M$ departure keys. Its finite evaluator accumulates integers without intermediate rounding, then rounds once. The same copy layer acquires the terminal pair. All structural masks, positions, local features, storage, and copying are charged.

A finite comparator implements equality acceptance in \eqref{eq:approx-gate-main}; restoring division produces an output with error at most $2^{-b}$. The construction covers empty rows, $h=0$, $b=0$, and numerator $B_i=Q$, whose binary representation requires $h+1$ bits. If the bank is presented as scalars rather than packed bits, explicit finite-grid ReLU extraction is included rather than treating scalar-to-bit conversion as free wiring. The resulting risk is
\begin{equation}\label{eq:finite-realization-main}
 R\le\min\left\{1,\frac{131072}{\sqrt M}+6433056\frac{k}{M}
      +1072176\bigl(2^{-h-1}+\delta\bigr)+2^{-b}\right\}.
\end{equation}
The predictor does not receive $k$ or $\delta$. This is a composition of a deterministic all-word construction with \eqref{eq:approx-risk-main}; it is not a training result.

\paragraph{What the small hybrid shape counts.}
Put $n=M+3$, $\ell_a=\lceil\log_2(n+1)\rceil$, and $p=\ell_a+h+b+32$. The hybrid has two attention layers and three heads in total, value width at most twelve, and shared pointwise Boolean/ReLU depth $O(p^2)$ and width $O(p)$. Its attention primitive is evaluated by paid finite integer arithmetic. Table~\ref{tab:costs} summarizes two complete, alternative implementations; their costs cannot be mixed to select the smallest entry in each row. All-token pointwise execution, holds, zero coefficients in the dense version, input, endpoint acquisition, and clearing are included. Appendix~\ref{sec:implementation-proof} supplies offline construction and the full ledgers.

\begin{table}[ht]
\centering
\small
\caption{Conservative finite implementation bounds. Each column is one complete representation. Peak storage uses the conservative scalar allocation and includes simultaneously live code, masks, input, and work arrays.}\label{tab:costs}
\begin{tabular}{@{}lll@{}}
\toprule
Resource & Sparse lists and circuits & Dense masks and circuits\\
\midrule
Online bit operations & $O(np^4)$ & $O(n^2p^2+np^6)$\\
Program and code bits & $O(p^4)$ & $O(p^5)$\\
Mask and position bits & $O(n\ell_a)$ & $O(n^2+n\ell_a)$\\
Total peak bits & $O(p^4+n\ell_a+np^2)$ & $O(p^5+n^2+n\ell_a+np^2)$\\
\bottomrule
\end{tabular}
\end{table}

A pure Boolean/ReLU network can instead unroll the entire bit-operation schedule. With its paid time $T$ and total live storage $S$, the resulting depth is $O(T+p)$, width $O(S+p)$, and sparse gates and holds $O((T+p)(S+p))$. These are not the smaller hybrid depth and width. The bank's $12(h+1)+2$ transmitted bits are likewise not the program's total memory.

\paragraph{The same-information counter comparison.}
A classical one-pass algorithm maintains four triples of exact $J,C,D$ counters and a two-bit sliding window, updating the half count only before $\lfloor M/2\rfloor$. It retains the final pair and can round once at the end to produce exactly the same bank. Its auxiliary storage is $O(\log M)$ bits and conservative counting work is $O(M\log M)$ bit operations, plus separately paid finite output division. Resident input costs $M+2$ additional bits; a read-once input assumption must be stated if that storage is omitted. This baseline achieves the exact-statistic rate and is computationally simpler than the attention realization. The comparison to count-only prediction isolates the added positional information, not an advantage of neural computation.

\subsection{Precision of the fixed bank}
Let $R^*(M,h)$ be the clean minimax actual-terminal TV risk when the only observation is the twelve separately nearest-rounded coordinates in \eqref{eq:direct-bank}, the exact endpoint, and source-independent public metadata. The decoder may have arbitrary computation, real-valued probability outputs, and independent randomization. The source class is the same critical-certificate class with arbitrary starts, but no source edits or numeric faults. Theorem~\ref{thm:bank-precision} proves
\begin{equation}\label{eq:bank-precision-main}
 \frac{1}{1600000}\max\{M^{-1/2},2^{-h}\}
 \le R^*(M,h)
 \le\min\{1,667160\max\{M^{-1/2},2^{-h}\}\}.
\end{equation}
Both lower terms hold on the stationary IID Bernoulli subclass. The precision proof constructs two separated parameters whose twelve mean curves share a rounding cell, proves concentration of the actual random bank in that cell, and couples the retained endpoint without asserting independence from the bank. A full-history two-point bound covers the complementary sample-size regime.

Consequently $h=\tfrac12\log_2 M+O(1)$ is the scale needed to retain a fixed multiple of the clean statistical rate for this particular bank. This is not a bound on all encoders with the same total number of bits. It is not a lower bound on the arithmetic precision, addresses, accumulator width, parameters, scores, or total storage of a Transformer. The clean lower also does not assert matching $\delta$ or $k/M$ terms for this new fault port. Equation~\eqref{eq:finite-realization-main} is the proved robust upper.


\section{Relation to prior work and scope}\label{sec:related}
\paragraph{Clean positional prediction.}
\citet[model (1), Theorem 24, and Section VI.A]{HanJanaWu2023} study stationary prediction at the actual terminal state, including suffix-based estimators and higher-order KL analysis. Theorem 24 gives the binary order-two KL result; the squared-TV result in Section VI.A concerns first-order chains with fixed state count at least three. These are direct antecedents, not independent-query results. Their clean full-trajectory experiment differs from the strict-edit and restricted-observation comparison here. Losses must also be distinguished: a squared-TV upper bound implies an unsquared-TV upper bound by Jensen, but a squared-TV lower bound does not yield its square root as an unsquared-TV lower bound. The spectral-gap-restricted results in that work use an assumption different from the supplied all-state minorization \eqref{eq:certificate}.

\paragraph{Unknown contamination.}
Budget adaptation is not new in itself. \citet{JainOrlitskyRavindrakumar2022} select among candidate estimates with monotone radii under a simultaneous-validity premise: every valid-index ball must contain the target. Separate expected-risk bounds do not establish that premise. Their selection lemma is pointwise, however, so the random terminal target is not by itself an obstruction: it could be applied realization by realization if all candidate predictions share the protected query and simultaneous accuracy is proved. \citet{BhattEtAl2022} explicitly construct confidence-based adaptation for a mean under strict retrospective replacements of IID observations. \citet{ChenGaoRen2016} give robust testing and estimation under IID Huber mixtures, while \citet{LiuGao2019} exhibit an adaptation obstruction involving both smoothness and contamination. None of these comparisons licenses a general claim that unknown contamination alone prevents adaptation.

\paragraph{Dependence and count-based inference.}
\citet[Theorem 2]{MaityMitra2025} analyze robust mean estimation for a stationary finite ergodic Markov chain as a tool for temporal-difference learning. Their procedure takes a corruption fraction and mixing time; Bernoulli corruption indicators replace rewards rather than edit the state trajectory. Its invariant-mean target and mixing-dependent confidence bound differ from the actual-terminal word experiment here. Robustness under dependence is therefore an established subject, not itself the proposed contribution.


\citet{Lecestre2023} estimate marginal or invariant distributions under dependence; their contamination corollary imposes independence conditions on the replacement channel. \citet{LeskelaDreveton2026} estimate transition matrices from independent parallel trajectories with arbitrarily corrupted whole trajectories. These are not the one-word terminal-target experiment. They also show why the present result should not be summarized as a universal failure of counts under corruption.

Clean transition-count factorization and Markov exchangeability remain valid. The report of \citet{DiaconisFreedman1978} and the finite-stack discussion of \citet{HalberstamHutchcroft2024} provide relevant background on count-compatible paths. For a clean order-two word, likelihood is its initial-pair mass times transition probabilities raised to triple counts. After a position-sensitive kernel $K$, an observed word's probability instead sums these clean probabilities weighted by $K(z\mid x)$. Clean factorization alone does not establish sufficiency in this altered experiment. The specific observation map, not an unrestricted claim about all compact summaries, is essential.

\paragraph{Limits of the conclusions.}
The lower bounds concern uniform unknown-budget comparison at the critical supplied certificate, with exact counts and fixed boundary windows. They are not full-history impossibility results, every-budget logarithmic minimax penalties, claims about arbitrary interior statistics, or noisy-bit guarantees. The positive results include arbitrary starts; the lower witnesses are stationary. The constants are deliberately conservative. The word example is not empirical validation, and no learned-predictor experiment is asserted.

\paragraph{Attention, counting, and Markov representation.}
Uniform selected averaging and positional routing are standard operations in RASP \citep{WeissGoldbergYahav2021}. Histogram constructions using uniform mixing also appear in \citet{BehrensEtAl2024}. \citet[Proposition 2.1]{EdelmanEtAl2024} give a two-layer bigram mechanism with previous-token copying; \citet{RajaramanEtAl2024} and \citet{EkboteEtAl2025} give shallow conditional $k$-gram representations under their respective positional and normalization models. These works preclude presenting shallow empirical Markov prediction as a new representation principle. The structural singleton masks used here are an explicitly paid different interface, not a bound for dense-causal-only attention. Their learning results and empirical studies address specified training settings; none is imported to claim that the present gate is learned. \citet{LepageMaryPicard2025} provide further empirical Markov-prediction context, comparing trained models with smoothed count estimates under different transition-matrix distributions.

\paragraph{Precision models.}
\citet[Proposition 1]{MerrillSabharwal2023} analyze rounding individual attention probabilities. Our finite evaluator instead sums integer indicators exactly and rounds the aggregate once. These operations have different precision requirements: $h$ describes the exported normalized bank, while addresses, accumulators, comparisons, and division use wider exact words. Balancing a statistical error with quantization error is standard estimation methodology; the assertion here is the matched upper and lower for this fixed bank. It is not an optimized-encoder or whole-network precision lower bound.

\paragraph{Earlier project results and scope of this consolidation.}
The prior finite-bank analysis \citep{BCJ001}, dated 2 November 2026, already proves real-error gate robustness and matching precision, numeric-noise, and source-edit rates for eight floored triple frequencies and four floored half-prefix frequencies, with a finite sequential evaluator. The present direct $J/C/D$ interface adds sequential clipping; its nearest-rounded bank has a different lower proof and a different canonical post-rounding fault port. The implementation makes local copying, Boolean/ReLU features, scalar-grid extraction, and hybrid shape explicit. These are interface and realization refinements of that earlier phenomenon. Separately rounded representations are not automatically equivalent after taking sums or differences, even when their unrounded counts determine each other.

\section{Discussion}\label{sec:discussion}
The main theorem identifies the price of a specified restriction on positional observations. Exact transition counts with short boundary windows cannot retain the known-budget scale for one map across all budgets at a fixed horizon. A proportional old-visit gate recovers the scale with a fixed positive fraction of the raw trajectory, or with a single computed bit after a whole-word scan. Its stability under approximate statistics makes a finite realization possible without exact recovery of the gate bit.

The complete argument relies on both sides of the information comparison: a gated reciprocal-occupation bound for the upper, and a legal histogram- and boundary-preserving transport for the lower. Neither a generic robust selector nor a generic constrained-risk inequality removes those experiment-specific obligations. The external comparisons given here are not an exhaustive novelty certification.

A same-information classical counter already achieves the statistical guarantee with logarithmic auxiliary storage. The attention construction therefore supplies a transparent realization, not a speed, memory, or statistical advantage over that counter. Learning the positional decision in a specified trainable architecture, or imposing dense causal masks and bounded logits throughout, would require separate assumptions and evidence. No such result is asserted here.

The exact supplemental bit in Section~\ref{sec:bit} is fault-free after construction. Transmission through noisy channels is a separate problem, treated in the project's companion analyses of
\href{https://github.com/AlexenderSokolov/transformer-icml/blob/670b7fc34ed864716124e9a56876d1baaaed6503/v10_research/resumption_20261103/RESEARCH_SUMMARY.zh.md}{known channels},
\href{https://github.com/AlexenderSokolov/transformer-icml/blob/670b7fc34ed864716124e9a56876d1baaaed6503/v10_research/channel_uncertainty_stage_20261103/RESEARCH_SUMMARY.zh.md}{uncertain channels}, and
\href{https://github.com/AlexenderSokolov/transformer-icml/blob/670b7fc34ed864716124e9a56876d1baaaed6503/v10_research/conditional_channel_20261103/CANDIDATE.md}{conditionally bounded noise}; it is not covered by the numeric-bank replacement model of Section~\ref{sec:finite-implementation}. In particular, a numeric coordinate perturbation, a bit error in an encoded message, and a corrupted arithmetic register are different experiments. The companion-work links and the claim-to-source ledger accompany this manuscript rather than adding unrelated coding theorems to its proof sequence.

\clearpage
\appendix
% MANUSCRIPT DRAFT: pending complete independent consolidated review.
% This file contains no external proof dependencies on project-private units.
\section{Finite-state estimates and upper bounds}
\label{sec:positive-proofs}

Throughout, a departure count excludes the terminal state.  For a chain
$(Y_i)_{i=0}^M$, write
\[
 J_u=\sum_{i=0}^{M-1}\mathbf1\{Y_i=u\},\qquad J=J_{Y_M}.
\]
For the binary order-two experiment, $Y_i=(X_i,X_{i+1})$,
$X_0,\ldots,X_{M+1}$ contains $M$ departures, and
$q_u=\theta_u(d(u))$ with $d(a,b)=1-a$.  Let $C_u$ count the
corresponding designated successors and put $C=C_{Y_M}$ and $q=q_{Y_M}$.
Primed counts are computed on the corrupted word.  Every bound below permits
an arbitrary initial law.  The all-state certificate, when invoked, is
$P^2(u,v)\ge\beta\Pi(v)$ for every pair of states, with $\beta>0$.
The final pair is protected, so clean and corrupted counts always select the
same actual query row.

\subsection{Universal occupation and clean prediction}

\begin{lemma}[Actual-terminal occupation]\label{lem:occupation}
For any homogeneous chain on $d$ states, any initial law, any $M\ge0$, and
any integer $w\ge1$,
\[
 \Pr(J_{Y_M}<w)\le\min\left\{1,\frac{d^2w}{M+1}\right\}.
\]
No irreducibility or stationary-start assumption is needed.
\end{lemma}
\begin{proof}
For each $t$ set $J(t)=\sum_{i<t}\mathbf1\{Y_i=Y_t\}$,
$b_t(v)=\Pr_v(J(t)<w)$, and $B_t=\max_v b_t(v)$.  If $t\ge\ell$,
the full count dominates the count in the final $\ell$-step suffix.
Conditioning on the suffix's initial state and using homogeneity gives
$B_t\le B_\ell$.  On every path, each state can occur at most $w$ times
with fewer than $w$ earlier occurrences.  Hence
$\sum_{t=0}^M\mathbf1\{J(t)<w\}\le dw$.  Consequently
\[
 (M+1)B_M\le\sum_{t=0}^M B_t
 \le\sum_v\sum_{t=0}^M b_t(v)\le d^2w.
\]
Average over the initial law and use the probability cap.
\end{proof}

\begin{lemma}[Binary clean empirical comparison]\label{lem:empirical}
For every homogeneous binary order-two source, every initial law, and
$M\ge64$,
\[
 \E\left[\mathbf1\{J>0\}\left|\frac CJ-q\right|\right]
 \le\frac{512}{\sqrt{M+1}}.
\]
The expression inside the indicator is evaluated only when $J>0$.
\end{lemma}
\begin{proof}
For each of the four rows $u$, generate an independent stack of $2M$
Bernoulli$(q_u)$ variables $\xi_{u,j}$, independently of the initial pair.
At the $j$th departure from $u$, use $\xi_{u,j}$ to choose whether the
successor is $d(u)$.  For every fixed finite path, the used stack prefixes
have exactly the product of its transition probabilities; thus this realizes
the original law, including boundary probabilities.  Write
$S_{u,j}=\sum_{i=1}^j(\xi_{u,i}-q_u)$.  Pathwise,
$C-Jq=S_{Y_M,J}$ on $J>0$.  The stacks are not asserted independent after
conditioning on $J$, $Y_M$, or any selection event.

For a centered variable in an interval of length one, its log moment
generating function has second derivative equal to a tilted variance,
which is at most $1/4$.  Integrating twice gives
$\E e^{\lambda V}\le e^{\lambda^2/8}$.  Thus
$e^{\lambda S_{u,j}-j\lambda^2/8}$ is a nonnegative supermartingale.
Stopping at its first level crossing before the deterministic time $2v$
gives, for $x>0$,
\[
 \Pr\left(\max_{j\le2v}S_{u,j}\ge x\right)
 \le e^{-\lambda x+2v\lambda^2/8}
 \le e^{-x^2/v},\qquad \lambda=2x/v.
\]
The stopping inequality here is finite: the expectation of the stopped
process is at most its initial value by summing conditional expected
increments through time $2v$.  Apply the same argument to $-S$ and take
a union over rows.  For $Z_v=\max_u\max_{j\le2v}|S_{u,j}|$,
\begin{equation}\label{eq:p-stack-tail}
 \Pr(Z_v\ge x)\le8e^{-x^2/v}.
\end{equation}
For dyadic integers $v\le M$, let $E_v=\{v\le J<2v\}$ and
$p_v=\min\{1,32v/(M+1)\}$.  Lemma~\ref{lem:occupation} gives
$\Pr(E_v)\le p_v$.  Without any independence assumption, tail integration
and $a=\sqrt{v\log(8/p_v)}$ yield
\[
 \E[Z_v\mathbf1_{E_v}]
 \le\int_0^\infty\min\{p_v,8e^{-x^2/v}\}\,dx
 \le p_va+\frac{p_vv}{2a}
 \le2p_v\sqrt{v\log(8/p_v)}.
\]
For the middle inequality, bound the Gaussian tail above $a$ by inserting
$x/a\ge1$ and integrating.  This layer contributes at most
$2p_v\sqrt{\log(8/p_v)}/\sqrt v$ to the desired risk.
Put $N=M+1$.  Summing layers with $v\ge N/32$ costs at most
\[
 \frac{2\sqrt{32\log8}}{1-2^{-1/2}}N^{-1/2}<56N^{-1/2}.
\]
If $v_0$ is the largest dyadic integer below $N/32$, then
$N/64\le v_0<N/32$ and
$\log(N/(4v_02^{-j}))\le3(j+1)$.  The remaining layers cost at most
\[
 \frac{64\sqrt3}{\sqrt{32N}}
 \sum_{j\ge0}2^{-j/2}\sqrt{j+1}
 \le\frac{64\sqrt3}{\sqrt{32N}}(1-2^{-1/2})^{-2}
 <236N^{-1/2}.
\]
Only a numerical series was extended; all stack maxima use at most $2M$
variables.  The layers partition $1\le J\le M$, including the last layer
when $M$ is dyadic.  The safe constant $512$ follows.
\end{proof}

\subsection{Return cycles and binary hitting scales}

\begin{lemma}[Finite-state return identity]\label{lem:kac}
In a finite irreducible chain with stationary law $p$, the first positive
return $T$ to $u$ has mean $\E_uT=1/p_u$.  If
$H_u=\max_v\E_vD_u$, where $D_u$ is the ordinary hitting time, zero at
$u$, then $1/p_u\le1+H_u$.
\end{lemma}
\begin{proof}
Irreducibility and finiteness give a fixed block length and a positive
minimum probability of hitting $u$ within a block, uniformly over starting
states.  Repeated conditioning gives geometric tails for hitting and return
times, so the following expectations are finite.  Let
$\nu(v)=\E_u\sum_{j=0}^{T-1}\mathbf1\{Y_j=v\}$.  First-step conditioning
inside this sum gives
\[
 (\nu P)(v)=\E_u\sum_{j=1}^{T}\mathbf1\{Y_j=v\}=\nu(v),
\]
since $Y_0=Y_T=u$.  A finite irreducible chain has a unique stationary
probability: for completeness, if $\eta$ is another one, the ratios
$\eta(v)/p_v$ are harmonic for the reversed transition matrix; a maximal
ratio propagates along every positive reversed edge and hence is constant.
Thus $\nu=(\E_uT)p$.  Exactly one visit to $u$ occurs before $T$, so
$\nu(u)=1$, proving the identity.  A positive return is one step followed
by an ordinary hit, proving $\E_uT\le1+H_u$.
\end{proof}

\begin{lemma}[Support and weighted hitting scale]\label{lem:hitting-scale}
A binary order-two source with a positive all-state certificate has either
full stationary support or a singleton absorbing diagonal support.  In the
full-support case, with $q=(a,b,e,\tau)$ and
$\overline H_u=\max\{H_u,1/\Pi(u)\}$,
\begin{equation}\label{eq:p-weighted-hitting}
 \sum_u\Pi(u)q_u\overline H_u\le\frac{40}{\sqrt\beta}.
\end{equation}
\end{lemma}
\begin{proof}
Stationarity of consecutive pair marginals gives
$\Pi=(y,x,x,w)$, and stationary flow into each diagonal gives
\begin{equation}\label{eq:p-flow}
 ya=xe,\qquad w\tau=xb.
\end{equation}
If $x>0$, the certificate entries
$P^2(10,01)=ea\ge\beta x$ and
$P^2(01,10)=b\tau\ge\beta x$ force $a,b,e,\tau>0$.
The flow identities then imply $y,w>0$.  If $x=0$, the stationary law
is on diagonals; each supported diagonal is absorbing by stationarity.
Two absorbing supported diagonals contradict the positive certificate.
This proves the classification.  Summing a certificate row gives
$\beta\le1$.  Full support and the certificate imply irreducibility.

Assume full support and put $\gamma=(a,1,1,\tau)\ge q$.
The two certificate entries just used, together with
\eqref{eq:p-flow}, also give
\begin{equation}\label{eq:p-square-entry}
 e^2\ge\beta y,\qquad b^2\ge\beta w.
\end{equation}
Here, for example, $ea\ge\beta x=\beta ya/e$ and $a,e>0$.
For target $00$, write $h_{01},h_{10},h_{11}$ for the mean hit times.
The first-step equations are
\[
 h_{01}=1+(1-b)h_{10}+bh_{11},\quad
 h_{10}=1+(1-e)h_{01},\quad h_{11}=\tau^{-1}+h_{10}.
\]
They give
\[
 h_{01}=\frac{b+2\tau}{e\tau},\quad
 h_{10}=h_{01}-\frac b\tau-1,\quad
 h_{11}=h_{01}+\frac{1-b}{\tau}-1,
\]
and therefore $H_{00}\le b/(e\tau)+2/e+1/\tau+1$.
For target $01$, the three hit times from $00,10,11$ are
$1/a$, $1+e/a$, and $1/\tau+1+e/a$, so
$H_{01}\le2/a+1/\tau+2$.
Exchanging symbols gives the other two bounds:
\[
 H_{10}\le1/a+2/\tau+2,\qquad
 H_{11}\le e/(ba)+2/b+1/a+1.
\]
The substitution exchanges $00$ with $11$, $01$ with $10$, and
$(a,b,e,\tau)$ with $(\tau,e,b,a)$; it does not assume $b=e$.
Using \eqref{eq:p-flow} and \eqref{eq:p-square-entry},
\begin{align*}
 \sqrt\beta\,yaH_{00}
 &\le\sqrt\beta\,(w+2x+ew/b+xe)
 \le\sqrt\beta\,(w+3x)+\sqrt w\le5,\\
 \sqrt\beta\,xH_{01}
 &\le2\sqrt\beta\,y/e+\sqrt\beta\,w/b+2\sqrt\beta\,x
 \le2\sqrt y+\sqrt w+2x\le5.
\end{align*}
The symbol-exchanged inequalities give the same bound for $10,11$.
Lemma~\ref{lem:kac} implies $\overline H_u\le H_u+1$; hence
$\Pi(u)\gamma_u\overline H_u\le6/\sqrt\beta\le10/\sqrt\beta$.
Sum over four states and use $q\le\gamma$.
\end{proof}

\subsection{First hits, endpoint probabilities, and joint tails}

\begin{lemma}[Target endpoint and occupation bounds]\label{lem:endpoint-tail}
For any $d$-state chain, possibly reducible, and $v\ne u$,
\begin{equation}\label{eq:p-first-hit}
 \Pr_v(D_u=t)\le\frac d{t+1},\qquad t\ge1.
\end{equation}
For an irreducible chain, fix a target with stationary mass $p>0$,
maximum ordinary hitting mean $H$, and $h=\max\{H,1/p\}$.
For every initial state $i$ and every integer $\ell\ge1$, with
$J_u(\ell)=\sum_{j=0}^{\ell-1}\mathbf1\{Y_j=u\}$,
\begin{align}
 \Pr_i(Y_\ell=u)&\le4dp(1+h/\ell),\label{eq:p-endpoint}\\
 \Pr_i(J_u(\ell)<\ell p/16)&\le2e^{-\ell/(8h)},\label{eq:p-lower-tail}\\
 \Pr_i(Y_\ell=u,J_u(\ell)>4\ell p+1)
 &\le16dp(1+h/\ell)e^{-\ell/(2h)}.\label{eq:p-joint-tail}
\end{align}
\end{lemma}
\begin{proof}
\emph{Reciprocal return coefficients and first hits.}
For any finite stochastic matrix $P$, let $T_j$ be the $j$th positive
return to $u$, allowing defective returns, and let
$F(z)=\E_u[z^{T_1};T_1<\infty]$.  Strong Markov at returns gives
$\Pr_u(T_j=s)=[z^s]F(z)^j$.  All strong-Markov uses here can be obtained
by conditioning on each possible stopping time and state and multiplying
the subsequent transition probabilities; no conditional independent-sample
claim is involved.  Let $P_{-u}$ delete row and column $u$.  Writing $P$
in target/rest block form, direct summation over the excursion before the
first return gives
$F(z)=zP_{uu}+z^2P_{u,-u}(I-zP_{-u})^{-1}P_{-u,u}$.
The block determinant identity therefore yields
\[
 1-F(z)=\frac{\det(I-zP)}{\det(I-zP_{-u})}.
\]
These are formal series, so no convergence issue arises.  The identity
$-\log\det(I-zA)=\sum_{s\ge1}\operatorname{tr}(A^s)z^s/s$
follows by differentiating the determinant, expanding $(I-zA)^{-1}$,
and integrating from zero.  Consequently
\begin{equation}\label{eq:p-return-trace}
 \sum_{j\ge1}\frac{\Pr_u(T_j=s)}j
 =\frac{\operatorname{tr}(P^s)-\operatorname{tr}(P_{-u}^s)}s
 \le\frac ds,\qquad s\ge1.
\end{equation}
Its left side is nonnegative; the upper bound uses nonnegative entries
of $P_{-u}^s$ and diagonal entries of $P^s$ at most one.
Now change only row $u$ to send $u$ deterministically to $v\ne u$.
In this auxiliary stochastic chain, the first positive return has law
$1+D_u$ under the original start $v$.  Retain only the first-return
term of \eqref{eq:p-return-trace} to prove \eqref{eq:p-first-hit}.
The auxiliary chain need not be an order-two source.

\emph{A target-specific endpoint bound.}
Return now to the irreducible chain.  Let $g(v)=\E_vD_u$.
First-step equations and Lemma~\ref{lem:kac} give
$(Pg-g)(v)=-1+\mathbf1\{v=u\}/p$.  Telescoping expectations shows
\begin{equation}\label{eq:p-expected-count}
 \E_vJ_u(s)=p\{s+\E_vg(Y_s)-g(v)\}\le p(s+H).
\end{equation}
A loop-rotation identity is
\[
 sP^s(u,u)=\sum_v\E_v[J_u(s)\mathbf1\{Y_s=v\}]
 \le dp(s+H).
\]
Indeed, for each $j<s$, the corresponding sum is
$\sum_vP^j(v,u)P^{s-j}(u,v)=P^s(u,u)$.
For $v\ne u$, split the first-hit convolution at $t=\lfloor s/2\rfloor$.
For early hits the return factor is at most $dp(1+2H/s)$.
For late hits \eqref{eq:p-first-hit} is at most $2d/s$, and
\eqref{eq:p-expected-count} from $u$ bounds the sum of return factors
by $p(\lceil s/2\rceil+H)$.  The resulting bound is at most
$dp(2+2/s+4H/s)\le4dp(1+H/s)$.  It also holds for $v=u$;
the early sum may be empty at $s=1$.  This proves the slightly stronger
version of \eqref{eq:p-endpoint} with $H$ instead of $h$.

\emph{Return moments and positive transforms.}
The family consisting of ordinary hitting times from nontarget states and
the positive return from $u$ is closed under conditioning on survival and
the current state.  Its means are at most $h$.  For each member $V$ and
integer $r\ge1$,
\begin{equation}\label{eq:p-return-moment}
 \E V^r\le r!h^{r-1}\E V.
\end{equation}
To prove this by induction, use
$V^r\le r\sum_{j=0}^{V-1}(V-j)^{r-1}$, condition each summand on
survival through time $j$, and apply the induction hypothesis and the
uniform mean bound $h$ to the residual law.  The sum of the survival
probabilities is $\E V$.
With $\lambda=1/(4h)$, expansion in nonnegative powers gives
\begin{equation}\label{eq:p-return-transforms}
 \E_i e^{\lambda D_u}\le\frac43,\qquad
 \E_u e^{\lambda T_1}\le e^{2\lambda/p},\qquad
 f:=\E_u e^{-\lambda T_1}\le e^{-3\lambda/(4p)}.
\end{equation}
For the last inequality use $e^{-x}\le1-x+x^2/2$ and
$\E T_1^2\le2h/p$; the first two follow from the geometric-series
bound on the moments.  If $K=\lceil\ell p/16\rceil$, fewer than $K$
visits in $0,\ldots,\ell-1$ means
$D_u+T_1+\cdots+T_{K-1}\ge\ell$.
These are fresh independent return cycles after the first hit, before any
horizon restriction.  Exponential Markov and \eqref{eq:p-return-transforms}
bound this probability by
$(4/3)e^{-\lambda\ell+2\lambda(K-1)/p}
\le(4/3)e^{-7\lambda\ell/8}\le2e^{-\ell/(8h)}$.
This includes $K=1$ and proves \eqref{eq:p-lower-tail}.

\emph{A joint tail by an exact change of path measure.}
Let $h_i=\E_i e^{-\lambda D_u}$, with $h_u=1$.
The inequality $e^{-x}\ge1-x$ gives $h_i\ge3/4$ and $f\ge3/4$.
Define
\[
 Q_{ij}=e^{-\lambda}P_{ij}h_j/h_i\quad(i\ne u),
 \qquad Q_{uj}=e^{-\lambda}P_{uj}h_j/f.
\]
The hitting and return equations show that every row sums to one.
This auxiliary matrix preserves directed support, so is irreducible.
Its ordinary hitting law is the original law tilted by
$e^{-\lambda D_u}/h_i$, whence $H_Q\le2H$.
Its positive return law is tilted by $e^{-\lambda T_1}/f$, and
\[
 \frac1{2p}\le\E_QT_1\le\frac2p.
\]
For the lower bound, use
$\E[T_1e^{-\lambda T_1}]\ge1/p-\lambda(2h/p)=1/(2p)$
and $f\le1$; the upper bound uses $f\ge3/4$.
Kac for $Q$ gives $p_Q\le2p$ and
$\max(H_Q,1/p_Q)\le2h$.  The endpoint bound already proved yields
$Q^\ell(i,u)\le16dp(1+h/\ell)$.
For every length-$\ell$ path from $i$ to $u$, telescoping the factors gives
\[
 Q(\text{path})=P(\text{path})e^{-\lambda\ell}
                  f^{-J_u(\ell)}/h_i.
\]
The count is exactly the number of target departures.  Thus on
$J_u(\ell)>4\ell p+1$, summing the reverse likelihood relation gives
\[
 \Pr_i(Y_\ell=u,J_u(\ell)>4\ell p+1)
 \le h_ie^{\lambda\ell}f^{4\ell p}Q^\ell(i,u)
 \le16dp(1+h/\ell)e^{-\ell/(2h)}.
\]
This proves the joint bound with its stationary-mass factor; no product
of an unconditional upper tail and an endpoint probability was taken.
\end{proof}

\subsection{Weighted imbalance, including unsupported starts}

\begin{lemma}[Half-window imbalance]\label{lem:imbalance}
For $M\ge64$, put $n=\lfloor M/2\rfloor$ and
$J_{1,u}=\sum_{i=0}^{n-1}\mathbf1\{Y_i=u\}$.
Every binary source with a positive all-state certificate, under every
initial law, satisfies
\[
 \E\left[q_{Y_M}\mathbf1\{J>0,J_{1,Y_M}<J/256\}\right]
 \le\frac{65536}{M\sqrt\beta}.
\]
\end{lemma}
\begin{proof}
We first prove a general irreducible-chain statement.  Let $m=M-n$,
$J_{2,u}=\sum_{i=n}^{M-1}\mathbf1\{Y_i=u\}$, and assign arbitrary
weights $0\le q_u\le1$.  Write
$W_u=p_uq_u\overline H_u$, where $p_u$ is stationary mass.
For every $0<c\le1/256$ we claim
\begin{equation}\label{eq:p-generic-imbalance}
 \E\left[q_{Y_M}\mathbf1\{J_{1,Y_M}<cJ\}\right]
 \le\frac{288d}{M}\sum_uW_u.
\end{equation}
The strict imbalance event is empty at $J=0$.
Fix $u$ with $np_u\ge1$.  If
$J_{1,u}\ge np_u/16$ and $J_{2,u}\le4mp_u+1$, then $m\le2n$ gives
$J_{1,u}/(J_{1,u}+J_{2,u})\ge1/145>1/256$.
Thus imbalance requires a first-window lower tail or a second-window
upper tail.  For the first event, condition on the whole history through
$Y_n$, use the uniform future endpoint bound
\eqref{eq:p-endpoint}, and then the first-window lower tail
\eqref{eq:p-lower-tail}.  For the second event condition only on $Y_n$
and use the joint bound \eqref{eq:p-joint-tail}.  Since $m\ge n$,
\[
 \Pr_i(Y_M=u,\ J_{1,u}<cJ_u)
 \le24dp_u(1+\overline H_u/n)e^{-n/(8\overline H_u)}.
\]
After multiplication by $q_u$, this is
$(24dW_u/n)(x+1)e^{-x/8}$, where $x=n/\overline H_u$.
Since the maximum of $(x+1)e^{-x/8}$ is
$8e^{-7/8}<4$, it is at most $96dW_u/n$.
If $np_u<1$, use only the endpoint bound at time $M$.
As $W_u\ge q_u$ and $p_uq_u<q_u/n$,
\[
 q_u\Pr_i(Y_M=u)\le4d(p_uq_u+W_u/M)\le8dW_u/n.
\]
Sum over $u$, average over initial states, and use $n\ge M/3$ to prove
\eqref{eq:p-generic-imbalance}.
For full-support binary sources, Lemma~\ref{lem:hitting-scale} now gives
$288\cdot4\cdot40/(M\sqrt\beta)=46080/(M\sqrt\beta)$.

It remains to treat singleton stationary support, without discarding
unsupported actual initial states.  Suppose $00$ is absorbing, so $a=0$,
and denote the other states by $A=01$, $B=10$, $C=11$.
An original path ending in this active set has never hit $00$.
Reflect the killing edge $B\to00$ to $B\to A$, producing the auxiliary
three-state chain
\[
 R(A,B)=1-b,\quad R(A,C)=b,\quad R(B,A)=1,
 \quad R(C,B)=\tau,\quad R(C,C)=1-\tau.
\]
Every original active path has probability at most its reflected-path
probability: the only changed factor is $1-e\le1$ on $B\to A$.
The endpoint and all active-state counts are unchanged.  Thus any
nonnegative functional of these paths is dominated without conditioning
on survival.
For $b,\tau>0$ this chain is irreducible, with
\begin{equation}\label{eq:p-reflected-weights}
 Z=2\tau+b,\quad p_A=p_B=\tau/Z,\quad p_C=b/Z,
 \qquad p_Ab\overline H_A\le1,\quad
 p_C\tau\overline H_C\le1.
\end{equation}
Indeed $H_A=1+1/\tau$ and $H_C=2/b$ by first-step equations;
$\overline H_A\le2+1/\tau$ and
$\overline H_C\le1+2/b$ give the last inequalities since $b,\tau\le1$.
Apply \eqref{eq:p-generic-imbalance} with weights $(b,0,\tau)$:
the $A,C$ contribution is at most $288\cdot3\cdot2/M=1728/M$.
If $b=0$ or $\tau=0$, approximate these two parameters by positive
ones in the reflected chain.  At fixed finite horizon, path probabilities
and the weights are continuous, and the bound is uniform, so the same
conclusion follows by taking a limit.  No minorization-preserving
approximation of the original source is being asserted.
The remaining target $B$ has weight $e$, and for an active start
\begin{equation}\label{eq:p-next-absorption}
 e\Pr(Y_M=B)=\Pr(D_{00}=M+1)\le\frac4{M+2},
\end{equation}
by \eqref{eq:p-first-hit}: the sole entry edge to $00$ is $B\to00$.
Drop its imbalance event.  An initial $00$ contributes zero; mixtures
follow by linearity.  The total is at most $1732/M$, below the stated
bound because $\beta\le1$.  Exchanging symbols handles absorbing $11$
with the same fixed designated-successor rule.
\end{proof}

\subsection{Reciprocal occupation with an old visit}

\begin{lemma}[Old-visit reciprocal count]\label{lem:old-inverse}
For any homogeneous $d$-state chain and initial law, let $M\ge2$,
$0\le m\le\lfloor M/2\rfloor$, and
\[
 O_u=\sum_{i=0}^{M-m-1}\mathbf1\{Y_i=u\},\qquad O=O_{Y_M}.
\]
For every $c>0$,
\begin{align*}
 \E\frac{\mathbf1\{O\ge cJ\}}{J+1}
 &\le\frac1M\left(3d^2+2d^3
     \sum_{s=m+1}^{\lfloor(M-1)/2\rfloor}\frac1s\right)\\
 &\le\frac{3d^2+2d^3}{M}
       \left(1+\log\frac M{m+1}\right),
\end{align*}
where an empty sum is zero.  In particular, the last numerator constant
is $176$ for the pair chain.  The statement includes $J=0$ and $m=0$.
\end{lemma}
\begin{proof}
Starting at $u$, the event
$\{Y_s=u,J_u(s)=j\}$ is exactly $\{T_j=s\}$ for $j\ge1$:
the initial visit is a departure, while the terminal visit is not.
Equation~\eqref{eq:p-return-trace}, which allowed defective returns, gives
\[
 K_u(s):=\E_u\frac{\mathbf1\{Y_s=u\}}{J_u(s)+1}
 =\sum_{j\ge1}\frac{\Pr_u(T_j=s)}{j+1}\le\frac ds.
\]
Fix an initial state $v\ne u$, and put
$a_v(t)=\Pr_v(D_u=t)$.  The event $Y_M=u,J_u=0$ means $D_u=M$,
so its contribution is at most $d/(M+1)\le d/M$ by
\eqref{eq:p-first-hit}; it is impossible if $v=u$.
On $J_u>0,O_u\ge cJ_u$, integrality forces $O_u\ge1$, hence the first
hit is at most $M-m-1$.  Discarding the rest of the gate and using
strong Markov gives the exact identity
\[
 \E_v\frac{\mathbf1\{Y_M=u,O_u\ge1\}}{J_u+1}
 =\sum_{t=1}^{M-m-1}a_v(t)K_u(M-t).
\]
There is no target departure before $t$, so the full count is exactly
the post-hit departure count, with no additional unit shift.
For $t\le M/2$, use $K_u(M-t)\le2d/M$ and
$\sum_ta_v(t)\le1$; the contribution is at most $2d/M$.
This also covers $v=u$, whose expression is just $K_u(M)$.
For $t>M/2$, use $a_v(t)\le2d/M$ and write $s=M-t$;
the contribution is at most
\[
 \frac{2d^2}{M}\sum_{s=m+1}^{\lfloor(M-1)/2\rfloor}\frac1s.
\]
Together with the zero-count term this is a per-target bound.
Average over the start and sum over $d$ terminal targets to obtain the
first inequality.  Finally,
$\sum_{s=a}^b1/s\le1/a+\log(b/a)\le1+\log(M/a)$ for $b\le M$,
with the empty case automatic, gives the stated logarithmic bound.
\end{proof}

\subsection{Protected edits and the two positional gates}

\begin{lemma}[Pathwise gate reduction]\label{lem:gate}
Let $D_u$ be a count of row $u$ over a fixed set of departure indices,
and at the protected row put $D=D_{Y_M}$.  On the corrupted word use
\[
 \widehat q=\begin{cases}
 C'/J',&J'>0\ \hbox{and }512D'\ge J',\\
 0,&\hbox{otherwise}.
 \end{cases}
\]
For every allowed corruption kernel and every initial law,
\begin{align}\label{eq:p-gate-reduction}
 R\le{}&\E\left[\mathbf1\{J>0\}|C/J-q|\right]
 +48k\E\frac{\mathbf1\{D\ge J/1024\}}{J+1}
 +\E[q\mathbf1\{J>0,D<J/256\}]\notag\\
 &+\Pr(J<6156k+1)+\Pr(J<3078k+1).
\end{align}
In particular this applies to a first-half count or to the old-prefix
count reconstructed by subtracting a suffix count from the total.
\end{lemma}
\begin{proof}
One edited symbol affects at most three departure triples; each changed
triple changes the histogram's $\ell^1$ norm by at most two.  Pair-count
windows affect at most two departures.  Taking a union over the edits,
including adjacent edits and indices straddling a split, gives the safe
bounds
\[
 |J'-J|,\ |C'-C|,\ |D'-D|\le\delta:=6k.
\]
When the old count is computed by subtraction, it is still exactly a
direct count on each word, so the same bound applies without adding
two separate perturbation bounds.  Endpoint protection keeps the selected
row fixed.
Let $G'=\{J'>0,512D'\ge J'\}$,
$H=\{D\ge J/1024\}$, and
$B=\{J>0,D<J/256\}$.  On $G'\cap H$ with $J>0$,
\[
 |C'/J'-q|\le|C/J-q|+\frac{8\delta}{J+1}.
\]
For $\delta=0$ the ratios coincide.  For $\delta>0$ and $J\ge2\delta$,
$J'\ge J/2$ and ratio subtraction gives a difference at most
$4\delta/J\le8\delta/(J+1)$.  For $J<2\delta$, use the unit bound on
the difference of two probabilities.
On $G'\cap H^c$,
$D+\delta\ge D'\ge(J-\delta)/512$ whereas $D<J/1024$,
so $J<1026\delta=6156k$.
On $(G')^c\cap B^c$ with $J>0$, if $J'=0$ then $J\le\delta$;
otherwise $D'<J'/512$ and $D\ge J/256$ give
$J/256-\delta<(J+\delta)/512$, hence $J<513\delta=3078k$.
If the gate rejects on $B$, its loss is exactly $q$.
The two displayed integer-threshold tails also cover $J=0$.
These cases give \eqref{eq:p-gate-reduction} pathwise and then in
expectation, with $8\delta=48k$.  At $k=0$ the false-status cases for
$J>0$ are impossible.  Equality in the gate is accepted throughout.
Since the comparison is pathwise, arbitrary permitted attack randomization
does not change it.
\end{proof}

\begin{proposition}[Finite suffix and half-gate upper bounds]
\label{prop:finite-upper}
Let $M\ge64$, $2\le r\le M+2$, and set
\[
 m=\min\{r-2,\lfloor M/2\rfloor\},\quad
 N=M-m,\quad L=1+\log\frac M{m+1}.
\]
From the histogram obtain $J'_u,C'_u$, and from the last $m+2$ symbols
obtain $T'_u=\sum_{i=M-m}^{M-1}\mathbf1\{Y'_i=u\}$.
Use the gate in Lemma~\ref{lem:gate} with $D'_u=O'_u=J'_u-T'_u$.
This $k,\beta$-blind estimator has all-start risk
\begin{equation}\label{eq:p-suffix-upper}
 R\le\min\left\{1,\frac{65536}{\sqrt M}
       +\frac{65536}{M\sqrt\beta}
       +(8448L+147744)\frac kM\right\}.
\end{equation}
Alternatively, use $D'_u=J'_{1,u}$, the count over the first
$\lfloor M/2\rfloor$ departures.  Its all-start risk is
\begin{equation}\label{eq:p-half-upper}
 R\le\min\left\{1,\frac{65536}{\sqrt M}
       +\frac{65536}{M\sqrt\beta}+2^{23}\frac kM\right\}.
\end{equation}
The latter estimator factors exactly through the histogram, protected
pair, and the one bit $\mathbf1\{512J'_{1,Y_M}\ge J'_{Y_M}\}$.
\end{proposition}
\begin{proof}
The last $m+2$ raw symbols reveal exactly $m$ departure pairs and the
terminal pair.  Thus $O'_u$ is precisely the count over indices
$0,\ldots,N-1$, even though their raw order is not observed.
In particular $r=2$ means $m=0$ and $O'=J'$; counting $r-1$ suffix
departures would incorrectly include the terminal pair.
Since $N\ge\lceil M/2\rceil$, the clean $O$ is at least $J_1$.
Lemma~\ref{lem:imbalance} therefore bounds the bias in
\eqref{eq:p-gate-reduction} by $65536/(M\sqrt\beta)$.
Lemma~\ref{lem:old-inverse}, with $d=4,c=1/1024$, bounds its reciprocal
term by $48k\cdot176L/M=8448Lk/M$.
Lemma~\ref{lem:occupation} bounds the two tails together by
\[
 \frac{16(6156+3078)k+32}{M+1}
 =\frac{147744k+32}{M+1}.
\]
Lemma~\ref{lem:empirical} and the remaining $32/(M+1)$ are bounded by
$65536/\sqrt M$.  The unit loss cap proves \eqref{eq:p-suffix-upper}.

For the half gate we require a different reciprocal estimate, valid for
any four-state chain and start:
\begin{equation}\label{eq:p-half-inverse}
 \E\frac{\mathbf1\{J_1\ge J/1024\}}{J+1}\le\frac{130944}{M}.
\end{equation}
To prove it, fix $u$ and condition on the entire first-half history,
whose terminal state is $v$ and whose count of $u$ is $j$.
Let $m_2=M-\lfloor M/2\rfloor$ be the remaining length.  Balance implies
$J_{2,u}\le1023j$.  In the suffix, the fixed-terminal event
$Y_M=u,J_{2,u}\le1023j$ is a subset of the actual-terminal low-occupation
event.  Lemma~\ref{lem:occupation} gives conditional probability at most
$16(1023j+1)/(m_2+1)$.  Since $(J_u+1)^{-1}\le(j+1)^{-1}$,
the conditional contribution is at most $16368/(m_2+1)$, including $j=0$.
Sum over four targets and use $m_2+1\ge M/2$ to prove
\eqref{eq:p-half-inverse}.  This argument uses a deterministic prefix,
not endpoint-conditioned independent samples; Lemma~\ref{lem:occupation}
is valid at every suffix length.
Apply Lemmas~\ref{lem:gate}, \ref{lem:imbalance}, and
\ref{lem:empirical} again.  The total corruption coefficient is
\[
 48\cdot130944+16(6156+3078)=6433056<2^{23},
\]
proving \eqref{eq:p-half-upper}.
The stated bit is exactly the gate decision on every corrupted word.
If $J'=0$ its value is one, but the decoder's separate $J'>0$ test still
returns zero.  Thus neither this boundary case nor threshold equality
requires another message.  All edits occur before the exact message is
formed; the bit channel itself is not being corrupted.

At $\beta=1/M$, \eqref{eq:p-suffix-upper} gives directly, for
$g_M(k)=\max\{M^{-1/2},k/M\}$,
\[
 R\le(131072+156192L)g_M(k)
 \le287264L g_M(k)<2^{19}L g_M(k).
\]
For $r=\max\{2,\lceil\alpha M\rceil\}$ with $0<\alpha\le1$,
$m+1=\min\{r-1,\lfloor M/2\rfloor+1\}\ge\alpha M/2$.
Indeed $r-1\ge\alpha M/2$ follows by separating $\alpha M\le2$
from $\alpha M>2$.  Hence $L\le1+\log(2/\alpha)$.
Likewise \eqref{eq:p-half-upper} becomes
$\min\{1,2^{17}/\sqrt M+2^{23}k/M\}$ at the critical certificate.
\end{proof}

\subsection{A count-only oracle when the budget is supplied}

\begin{proposition}[Critical known-budget upper]\label{prop:count-oracle}
At supplied certificate $\beta=1/M$, a predictor using only the exact
triple histogram, the protected pair, and the public budget $k$ satisfies,
for every source and initial law,
\[
 R\le191488\max\{M^{-1/2},k/M\},\qquad M\ge64.
\]
For $k\ge1$ its designated probability is
\[
 \widehat q_k=\frac{(C'-6k)_+}{J'+6k+2};
\]
for $k=0$ use $(C'+1)/(J'+2)$.  This budget-dependent estimator is used
only for the oracle comparison, not for the unknown-budget results.
\end{proposition}
\begin{proof}
\emph{A weighted occupation estimate.}
For an irreducible $d$-state chain with target masses $p_u$, let
$h_u=\max\{H_u,1/p_u\}$, $0\le q_u\le1$, and $W_u=p_uq_uh_u$.
For every $A>0$ and every initial law we first prove
\begin{equation}\label{eq:p-oracle-weighted-general}
 \E\min\left\{q_{Y_M},\frac A{J+2}\right\}
 \le384d^2\frac AM+\frac{192d}{M}\sum_uW_u.
\end{equation}
Fix a target $u$.  If $h_u>M$, the endpoint bound gives
$q_u\Pr(Y_M=u)\le4dp_uq_u(1+h_u/M)\le8dW_u/M$.
If $h_u\le M$, on $J_u\ge Mp_u/48$ the minimum is at most
$48A/(Mp_u)$, and $\Pr(Y_M=u)\le8dp_u$; this costs at most
$384dA/M$.  For the complementary event, put $n=\lfloor M/2\rfloor$
and $m_2=M-n$.  Since $n\ge M/3$,
$J_u<Mp_u/48$ implies $J_{1,u}<np_u/16$.
Conditioning on the complete prefix through $Y_n$, use the uniform
endpoint bound on the remaining $m_2$ transitions, and then the
unconditional lower tail on the prefix.  The weighted contribution is
at most
\[
 8dp_uq_u(1+2h_u/M)e^{-M/(24h_u)}
 =\frac{8dW_u}{M}(x+2)e^{-x/24}
 \le\frac{192dW_u}{M},\qquad x=M/h_u.
\]
The last step uses
$\sup_{x\ge0}(x+2)e^{-x/24}=24e^{-11/12}<24$.
Summing these two cases over targets proves
\eqref{eq:p-oracle-weighted-general}.
For full-support binary sources, Lemma~\ref{lem:hitting-scale} yields
\begin{equation}\label{eq:p-oracle-weighted-binary}
 \E\min\left\{q_{Y_M},\frac A{J+2}\right\}
 \le6144\frac AM+\frac{30720}{M\sqrt\beta}.
\end{equation}

For singleton support $00$, the $00$ contribution vanishes.  Use the
reflected chain from the proof of Lemma~\ref{lem:imbalance}.  On paths
ending active, original path weights are bounded above by reflected
path weights, with exactly the same terminal state and departure counts.
Thus the nonnegative minimum functional for targets $A,C$, with weights
$(b,0,\tau)$, is dominated pathwise in the same sum over paths.
Equation~\eqref{eq:p-reflected-weights} and
\eqref{eq:p-oracle-weighted-general}, with $d=3$, bound this part by
$3456A/M+1152/M$.  The original target $B$ contributes at most
$e\Pr(Y_M=B)\le4/(M+2)$ by \eqref{eq:p-next-absorption}.
An initial $00$ contributes zero.  For boundary $b$ or $\tau$, approximate
only the reflected chain by positive parameters: there are finitely many
length-$M$ paths, and each summand is continuous in probabilities and
weights, including its minimum.  The uniform bound therefore passes to
the limit.  Symbol exchange handles singleton $11$.
Since $\beta\le1$, the larger bound
\eqref{eq:p-oracle-weighted-binary} covers both singleton cases and
arbitrary initial mixtures as well.

\emph{A pathwise lower envelope.}
For $k\ge1$ put $B_k=6k$, $D=J+2$, $D'=J'+B_k+2$, and
$c'=(C'-B_k)_+$.  The protected-edit perturbation bounds give
\[
 D'\ge D,\quad0\le c'\le C,\quad
 0\le C-c'\le2B_k,\quad0\le D'-D\le2B_k.
\]
Thus $0\le\widehat q_k\le C/(J+2)$, and
\[
 \frac CD-\frac{c'}{D'}
 =\frac{C-c'}D+\frac{c'}{D'}\frac{D'-D}D
 \le\frac{4B_k}D.
\]
For $J>0$, putting $p=C/J$ gives
\[
 0\le p-\widehat q_k
 \le\min\left\{p,\frac{4B_k+2}{J+2}\right\}.
\]
Since $x\mapsto\min\{x,c\}$ is $1$-Lipschitz on nonnegative reals,
with $A=24k+2$ the triangle inequality implies
\begin{equation}\label{eq:p-oracle-pathwise}
 |\widehat q_k-q|
 \le2\mathbf1\{J>0\}|C/J-q|
       +\min\left\{q,\frac A{J+2}\right\}.
\end{equation}
When $J=0$, $C=0$ and $\widehat q_k=0$; since $A\ge2$, the last
term equals $q$, so the inequality still holds without defining $C/J$.
This is valid for every permitted edit realization.
Using Lemma~\ref{lem:empirical},
\eqref{eq:p-oracle-weighted-binary}, $\beta=1/M$, and $A\le26k$ gives
\[
 R\le\frac{31744}{\sqrt M}+159744\frac kM
 \le191488\max\{M^{-1/2},k/M\},\qquad k\ge1.
\]

For $k=0$ the word is unedited.  On $J>0$, the add-one probability
is within $1/(J+2)$ of $C/J$; on $J=0$ pay unit loss.
Lemma~\ref{lem:occupation} gives $\Pr(J=0)\le16/(M+1)$.
For completeness, integer-threshold layer cake gives
\[
 \E\frac1{J+2}
 \le\int_0^1\min\left\{1,\frac{16}{(M+1)t}\right\}\,dt
 \le\frac8{\sqrt{M+1}}.
\]
Indeed a nonempty event $J<t^{-1}-2$ uses threshold
$\lceil t^{-1}-2\rceil\le t^{-1}$; for the last inequality use
$\min\{1,z\}\le\sqrt z$ and integrate $t^{-1/2}$.
Together with Lemma~\ref{lem:empirical} this is at most
$536/\sqrt M$, below the asserted constant.  No positional message or
raw-window input is used in either budget branch.
\end{proof}

\section{Observation separation and lower bounds}
\label{sec:negative-proofs}

Throughout this section a word has symbols $X_0,\ldots,X_{M+1}$,
pair states $Y_i=(X_i,X_{i+1})$, and departures $0\le i<M$.
Write $J_{ab}=\sum_{i<M}\ind\{Y_i=ab\}$,
$J_{\rm cross}=J_{01}+J_{10}$, and $N_{abc}=H_{abc}(X)$.
Thus $Y_M$ is a terminal pair, not a departure. All risks below include
independent predictor randomness. The observation $\mathcal T_r$ consists
of the eight triple counts and the actual first and last $r$ symbols.

\subsection{An exact occupation mixture}

For $0<a,t<1$, let $Q_a$ be the stationary law with pair transition matrix,
in the order $00,01,10,11$,
\[
 P_a=\begin{pmatrix}
 1-a&a&0&0\\
 0&0&1/2&1/2\\
 1/2&1/2&0&0\\
 0&0&t&1-t
 \end{pmatrix},\qquad
 \Pi_a=\frac{(t/a,2t,2t,1)}{1+t/a+4t}.
\]
In particular the stationary initial factor is
$\Pi_a(11)=a/[t+(1+4t)a]$.

\begin{lemma}[Occupation-mixture domination]
\label{lem:mixture}
Put $F=\{Y_0=11,Y_M=00\}$ and $U=N_{001}$. For every integer $j$,
every integer $u_0\ge1$, every interval $I\subset(0,1)$, and every
additional event $G$ (possibly depending on $a$),
\[
 \int_I Q_a(F,J_{00}=j,U\ge u_0,G)\,\frac{da}{a}
 \le \frac{t}{u_0}.
\]
\end{lemma}
\begin{proof}
Delete every transition departing from $00$, both loops and exits.
For $J_{00}=j$, exactly $L=M-j$ transitions remain. Their last transition
is $10\to00$. Before that transition the retained path is on
$B=01,C=10,D=11$. A retained $C\to00$ followed by an erased exit
$00\to B$ is represented by a marked $C\to B$ of weight $1/2$;
a direct $C\to B$ is unmarked and also has weight $1/2$. The erased
exit's factor $a$ is restored separately. All other retained transitions
keep their probabilities. The marked transition matrix is therefore
\[
 K(z)=\begin{pmatrix}
 0&1/2&1/2\\ (1+z)/2&0&0\\ 0&t&1-t
 \end{pmatrix}.
\]
For $L\ge1$, define
\[
 h_{L,u}=\tfrac12[z^u](K(z)^{L-1})_{D,C},
\]
and set $h_{L,u}=0$ for $L\le0$ or $u<0$. The final transition
$C\to00$ has weight $1/2$ and is not marked. A skeleton with $u$ marks
has $u$ completed visits to $00$ and one terminal visit. Distributing
its $j-u$ loops among these $u+1$ holding slots gives exactly
$\binom ju$ weak compositions. Each composition and marked skeleton
recover one word, and deletion recovers both uniquely. Consequently,
for $0\le u\le j$,
\begin{equation}
 Q_a(F,J_{00}=j,U=u)
 =\Pi_a(11)\binom ju a^u(1-a)^{j-u}h_{M-j,u}.
 \label{eq:marked-word-identity}
\end{equation}
Infeasible cases have zero mass. The terminal holding slot is present
even though $Y_M$ is not a departure; there is no initial $00$ slot.

For $u\ge1$, the beta integral gives
$\binom ju\mathrm B(u,j-u+1)=1/u$. Thus, if
$A_{j,u}$ has the $\operatorname{Beta}(u,j-u+1)$ law,
\begin{align}
 \int_0^1 Q_a(F,J_{00}=j,U=u)\,\frac{da}{a}
 &=\frac{h_{M-j,u}}{u}\E\Pi_{A_{j,u}}(11)
 \le\frac{h_{M-j,u}}{u}. \label{eq:stationary-beta-factor}
\end{align}
The expectation is part of the equality: stationarity has not been
cancelled. For $u=0$ the corresponding factor is instead
$h_{M-j,0}\int_0^1(1-a)^j/[t+(1+4t)a]\,da$, and no $1/u$ formula
is used. This is why $u_0>0$ is required.

To sum over marks at the exact outside horizon, put $R=K(1)$.
This matrix is stochastic. Started at $D$, let $b_n,c_n,d_n$ be its
probabilities of $B,C,D$. Then
\[
 c_0=0,\quad c_1=t,\quad b_n=c_{n-1}\ (n\ge1),\qquad
 c_n=\tfrac12 c_{n-2}+td_{n-1}\ (n\ge2).
\]
Induction gives $c_n\le2t$. Hence, for each $L\ge1$,
\[
 \sum_{u\ge0}h_{L,u}=\tfrac12 c_{L-1}\le t.
\]
First drop $G$ and enlarge $I$ to $(0,1)$ in the nonnegative integral.
Then sum \eqref{eq:stationary-beta-factor} over $u\ge u_0$ to obtain
\[
 \int_I Q_a(F,J_{00}=j,U\ge u_0,G)\,\frac{da}{a}
 \le\sum_{u\ge u_0}\frac{h_{M-j,u}}u
 \le\frac1{u_0}\sum_{u\ge0}h_{M-j,u}\le\frac t{u_0}.
\]
No conditioning on a prescribed number of outside marks or replacement
of the horizon $M-j$ is involved.
\end{proof}

\subsection{Relocating loops while preserving sampled windows}

\begin{lemma}[Interior loop relocation]
\label{lem:relocation}
Suppose $Y_0=11,Y_M=00$, $2\le r$, and $2r\le M+2$.
Put $l=r-2$, $h=M-r+2$. If the interior vertices
$Y_l,\ldots,Y_h$ contain a visit to $11$ strictly before a visit to $00$,
there is a deterministic transformation $C_r$ preserving the complete
triple histogram and the actual first and last $r$ symbols, such that,
with $j=J_{00}$,
\begin{equation}
 d_H\bigl(C_r(X),1^{M+2-j}0^j\bigr)
 \le3J_{\rm cross}+4r-5.
 \label{eq:relocation-distance}
\end{equation}
If at least one of the two movable diagonal-loop types is absent, the
same conclusion holds without the ordered-visit condition.
\end{lemma}
\begin{proof}
Freeze departures $i<l$ and $i\ge h$. These are $2r-4$ edges and
spell exactly the sampled boundary symbols. From the interior path
remove only the $111$ and $000$ loops, in numbers $A_1$ and $A_0$,
retaining every other edge in order. Let $S_0,\ldots,S_L$ be its vertices.
Deleting self-loops neither removes a distinct-state visit nor reverses
its order. Choose the first retained $11$ and the first subsequent
retained $00$; insert all $A_1$ loops at the former and all $A_0$ at
the latter. The choice is deterministic. Deleting or inserting a loop
preserves its attachment pair. Reinserting exactly the removed number
of edges preserves the absolute suffix attachment index $h$, so the
frozen edge labels and all sampled boundary symbols are unchanged.
Every triple label is retained or reinserted once, proving histogram
preservation. Unfrozen boundary-crossing triples need not keep their
individual positions.

Here is the required positional count argument. A binary word of length
$N$ with a distinguished block of $A_1$ ones before a distinguished
block of $A_0$ zeros has $o$ other ones and $z$ other zeros. If
$j=A_0+q$, $q\ge0$, then
\begin{equation}
 d_H(\text{word},1^{N-j}0^j)
 \le\max\{z+q,\,2o+z-q\}. \label{eq:two-block-distance}
\end{equation}
Indeed, write $z_-$ for zeros before the target cut. The exact mismatch
count is $2z_-+j-(A_0+z)$. If a distinguished one lies after the cut,
all distinguished zeros lie after the cut and $z_-\le z$, giving
$z+q$. Otherwise every one after the cut is among the $o$ other ones;
writing $o_+$ for their number, the other exact identity is
$2o_++(A_0+z)-j$, at most $2o+z-q$. This proof does not presume any
order among the remaining symbols.

Put $u=N_{001}$, $v=N_{011}$, $p=N_{010}$, $w=N_{101}$, and
$d=J_{\rm cross}$. Endpoint flow gives
\[
 N_{100}=u+1,\qquad N_{110}=v+1,\qquad
 d=p+v+u+1+w.
\]
Let $f_1,f_0$ be the frozen $111,000$ loop counts, so
$f_1+f_0\le2r-4$. The initial pair contributes two other ones.
Counting the appended symbol of every remaining edge gives exactly
\begin{align*}
 o&=2+u+v+w+f_1=d+1-p+f_1,\\
 z&=p+u+v+2+f_0=d+1-w+f_0,\\
 q&=j-A_0=u+f_0\ge0.
\end{align*}
Since $u\le d-1$ and $d\ge1$,
\begin{align*}
 z+q&=d+1-w+u+2f_0\le2d+2f_0,\\
 2o+z-q&=3d+3-2p-w-u+2f_1\le3d+3+2f_1.
\end{align*}
Each is at most $3d+4r-5$, proving \eqref{eq:relocation-distance}.
If only one movable loop type occurs, put its block at any retained
visit to its state, which exists because the loops were deleted there.
The same two-block argument treats the absent block as empty; if both
are absent use the identity transformation.

Finally, $j$ is the sampled word's actual departure count. The comparison
word has instead
\begin{equation}
 J_{00}(1^{M+2-j}0^j)=j-2\quad(2\le j\le M).
 \label{eq:null-suffix-shift}
\end{equation}
For this range its final pair and the transformed word's final pair are
both $00$, so all their Hamming differences are in the editable first
$M$ coordinates.
\end{proof}

\subsection{A finite matched-observation lower bound}

\begin{proposition}[Finite critical lower bound]
\label{prop:finite-lower}
Let
\[
 c=2^{-16},\quad A=2^{12},\quad t=\frac{ck}{M},\quad
 u_0=\lfloor ck/4\rfloor,\quad T=\lfloor M/8\rfloor,\quad
 L_-=\log\frac1{4At}.
\]
Assume integer $M,k,r$ satisfy
\begin{equation}
 \begin{gathered}
 M\ge64,\quad 2\le r\le k/128,\quad k\le M/8,\quad ck\ge1024,\\
 At<1/4,\qquad \frac M{k^2}\le\frac{Ac^2}{2^{20}}.
 \end{gathered}\label{eq:finite-lower-domain}
\end{equation}
For any fixed predictor based on $\mathcal T_r$, let $S$ be its supremum
clean risk over the stationary laws $Q_a$, $At\le a\le1/4$.
Every such source satisfies the all-state certificate with supplied
$\beta=1/M$. There is a stationary source with the same supplied
certificate and one source-independent strict-$k$ corruption kernel
for which
\begin{equation}
 R_{\rm null}\ge\frac{u_0}{8Mt}
       \left[\left(\frac t{16}-S\right)L_-\right]_+.
 \label{eq:finite-lower}
\end{equation}
The kernel protects the final pair on every input and every output.
\end{proposition}
\begin{proof}
We give the source, probability, and kernel calculations separately.

\paragraph{Source legality and all sixteen certificate entries.}
Throughout $t\le c/8$, $a\ge At\ge t$, and $a\le1/4$. Set
$D=1+t/a+4t$. Then
$1\le D\le1+1/A+c/2<4/3<2$.
The flow equalities
$\Pi_a(00)a=\Pi_a(10)/2$ and
$\Pi_a(11)t=\Pi_a(01)/2$, together with
$\Pi_a(01)=\Pi_a(10)$, verify stationarity directly. Multiplication gives
\begin{equation}
 P_a^2=\begin{pmatrix}
 (1-a)^2&(1-a)a&a/2&a/2\\
 1/4&1/4&t/2&(1-t)/2\\
 (1-a)/2&a/2&1/4&1/4\\
 t/2&t/2&(1-t)t&(1-t)^2
 \end{pmatrix}.
 \label{eq:all-sixteen}
\end{equation}
Every displayed entry is at least $t/2$: use $a\ge t$ and
$a,t\le1/4$ for the products and all other entries are immediate.
Since $t/2=ck/(2M)\ge1/M$, every starting and destination state obeys
$P_a^2(u,v)\ge M^{-1}\Pi_a(v)$. Also $a\ge At\ge2/M$, so
$a-\beta\ge a/2$.

\paragraph{Common-atom covariance and coupling bounds.}
The four paths to $11$ of length three, from the states in order, have
probabilities
\[
 (1-a)a/2,\qquad (1-t)^2/2,\qquad a/4,\qquad(1-t)^3.
\]
These arise respectively from paths $00,00,01,11$;
$01,11,11,11$; $10,00,01,11$; and $11,11,11,11$.
Each is at least $\lambda=a/4$, so
$P_a^3(x,\cdot)\ge\lambda\delta_{11}$ for every $x$.
Split each three-step transition into this common atom and a residual
probability law. Shared independent Bernoulli trials of probability
$\lambda$ couple any two chains at the first common atom. Fill in each
three-step block by its conditional bridge; this preserves the full
one-step path law. After meeting, use identical transitions. Thus the
contraction coefficient at lag $h$ is bounded by
\[
 \alpha_h=(1-\lambda)^{\lfloor h/3\rfloor}.
\]

For a stationary chain $W$, a departure-state-set count with stationary
mass $p$, or a count of a specified directed edge with stationary
probability $p$, satisfies
\begin{equation}
 \operatorname{Var}(V)\le\frac{7Mp}{\lambda}.
 \label{eq:common-atom-variance}
\end{equation}
For an edge indicator $I_i$ and $h\ge1$, condition on the state after
edge $i$. The conditional probability of edge $i+h$ is a $[0,1]$-valued
function after $h-1$ transitions, whose deviation from its stationary
value is at most $\alpha_{h-1}$. Therefore
$|\operatorname{Cov}(I_i,I_{i+h})|\le p\alpha_{h-1}$.
For a state-set indicator use $\alpha_h\le\alpha_{h-1}$ instead.
Since $\sum_{h\ge1}\alpha_{h-1}=3/\lambda$, expansion of the variance
gives $Mp(1+6/\lambda)\le7Mp/\lambda$, as claimed.

Couple a chain $Z$ started at $00$ to stationary $W$ by the same block
construction. Its meeting boundary $\sigma$ has
\[
 \E\sigma\le3/\lambda=12/a.
\]
Every departure-state or edge count in the two paths differs by at most
$\sigma$ pathwise: after boundary $\sigma$ the paths agree, and there
are at most $\sigma$ earlier departures. No independence of $\sigma$
and the stationary counts is needed.

\paragraph{Stationary reversal and count corrections.}
Let $\mathcal R(ab)=ba$. The stationary triple identity
\[
 \Pi_a(ab)P_a(ab,bc)=\Pi_a(cb)P_a(cb,ba)
\]
holds for every triple: the only nonpalindromic cases are $001,100$
and $011,110$, given by the two flow identities above. It implies
$\mathcal R P_a^*\mathcal R=P_a$, where $P_a^*$ is the stationary
reverse pair kernel. Reversing the symbols of a stationary word thus
preserves its law. In particular, with
$F=\{Y_0=11,Y_M=00\}$ and $E=\{Z_M=11\}$,
\[
 Q_a(F)=\Pi_a(00)P_{00}(E),
\]
and conditional on $F$, the reversed pair path
$Z_i=\mathcal RY_{M-i}$ has law $P_{00}(\cdot\mid E)$.
For $M\ge64$, $\lfloor M/3\rfloor\ge M/4$, and therefore
\[
 \alpha_M\le e^{-aM/16}\le\frac{16}{aM}
 \le\frac{16}{Ack}\le\frac14.
\]
Since $\Pi_a(11)=1/D\ge3/4$, we have $P_{00}(E)\ge1/2$.

Put $V(W)=N_{001}(W)$ and $x=ck$. Its stationary mean is
$Mt/D\in[x/2,x]$. Consider the four conditions
\begin{equation}
 \sigma\le x/16,\quad |V(W)-Mt/D|\le x/8,\quad
 J_{\rm cross}(W)\le k/32,\quad J_{00}(W)\le M/16.
 \label{eq:preliminary-good}
\end{equation}
Their failure probabilities, respectively, are at most
\begin{equation}
 \frac{192M}{Ac^2k^2},\quad
 \frac{1792M}{Ac^2k^2},\quad
 \frac{458752M}{Ak^2},\quad \frac{16}{A}.
 \label{eq:four-failures}
\end{equation}
The first follows from $\E\sigma\le12/a$ and $a\ge Ack/M$.
For the second, \eqref{eq:common-atom-variance} with $p=t/D$ gives
$\operatorname{Var}V\le28Mt/a$; divide by $(ck/8)^2$.
For the third, the mean is $4Mt/D\le4ck\le k/64$, and
$\operatorname{Var}J_{\rm cross}\le112Mt/a\le112M/A$;
divide by $(k/64)^2$. For the last use
$\Pi_a(00)=t/(aD)\le1/A$ and Markov's inequality.
Conditioning the joint coupling on $E$ costs at most two. Under
\eqref{eq:finite-lower-domain}, twice the sum in
\eqref{eq:four-failures} is bounded by
\[
 2\left(\frac{1984}{2^{20}}
       +\frac{458752c^2}{2^{20}}+\frac{16}{A}\right)<\frac1{16}.
\]

On $F$, the exact reversal identities are
\begin{equation}
 \begin{split}
 J_{00}(Z)&=J_{00}(Y)+1,\\
 J_{\rm cross}(Z)&=J_{\rm cross}(Y),\\
 N_{001}(Z)&=N_{100}(Y)=N_{001}(Y)+1.
 \end{split}\label{eq:reverse-count-shifts}
\end{equation}
The first compares arrival indices $1,\ldots,M$ with departure indices
$0,\ldots,M-1$; the second loses no mixed endpoint; the third combines
triple reversal with entry-minus-exit flow at $00$.
Thus \eqref{eq:preliminary-good}, the pathwise coupling bound, and
$x\ge1024$ imply for the original path
\begin{align*}
 U&\ge x/2-x/8-x/16-1\ge x/4\ge u_0,\\
 J_{\rm cross}&\le k/32+ck/16\le k/16,\\
 2\le u_0&\le J_{00}\le M/16+ck/16\le M/8.
\end{align*}
The integer upper count bound is $J_{00}\le T$.

\paragraph{Ordered visits in the interior.}
Let $l=r-2$ and $h=M-r+2$. Since $r\le M/1024$,
$h-l=M-2r+4\ge M/2$. Write $B$ for a visit to $11$ before a later
visit to $00$ within $Y_l,\ldots,Y_h$. Because $M-h=l$ and $M-l=h$,
reversal sends this event precisely to a visit to $00$ before a later
visit to $11$ within $Z_l,\ldots,Z_h$.

The ordinary expected hitting times of $00$ and $11$, zero at their
targets and listed in state order, are
\begin{align*}
 h_{00}&=(0,\,4+1/t,\,3+1/(2t),\,3+3/(2t)),\\
 h_{11}&=(3+3/(2a),\,3+1/(2a),\,4+1/a,\,0).
\end{align*}
For completeness, the first vector follows by solving
$h_{01}=1+(h_{10}+h_{11})/2$,
$h_{10}=1+h_{01}/2$, and
$h_{11}=1+th_{10}+(1-t)h_{11}$, with $h_{00}=0$.
For the second solve
$h_{00}=1+(1-a)h_{00}+ah_{01}$,
$h_{01}=1+h_{10}/2$, and
$h_{10}=1+(h_{00}+h_{01})/2$, with $h_{11}=0$.
On $a,t\le1/4$ the two maxima are
$3+3/(2t)$ and $3+3/(2a)$. From any state at time $l$, hitting $00$
and then $11$ therefore has expected additional time at most
\[
 6+\frac{3}{2t}+\frac{3}{2a}\le\frac6t.
\]
This uses the strong Markov property at the first hit: partition by its
finite time and state, apply the usual Markov property on each event,
and sum. Distinct targets make the ordering strict. Markov's inequality
over $h-l\ge M/2$ bounds the unconditioned failure probability by
$12/(tM)$. Dividing by $P_{00}(E)\ge1/2$ gives
\[
 Q_a(B^c\mid F)\le\frac{24}{ck}\le\frac3{128}.
\]
This estimate is uniform in the time-$l$ state; it assumes no
independence from the final state or the counts.

Let $G$ be the intersection of $B$ with
$U\ge u_0$, $J_{\rm cross}\le k/16$, and $2\le J_{00}\le T$.
The preceding count and hitting bounds give
$Q_a(G\mid F)\ge1-1/16-3/128>1/2$. Retaining the rare endpoint mass,
\begin{equation}
 Q_a(F\cap G)\ge\Pi_a(00)/4,\qquad
 (a-\beta)Q_a(F\cap G)\ge t/16.
 \label{eq:good-event-mass}
\end{equation}
Indeed $P_{00}(E)\ge1/2$, $a-\beta\ge a/2$, and
$a\Pi_a(00)=t/D\ge t/2$.

\paragraph{The null and its full word masses.}
Take a stationary first-order binary source with
$T_*(1,0)=q_*=1/(2M)$ and $T_*(0,1)=\beta=1/M$.
Its symbol masses are $\pi_*(0)=1/3$, $\pi_*(1)=2/3$.
Its stationary pair law is $\Pi_*(s,z)=\pi_*(s)T_*(s,z)$, and
for all starting and destination pair states,
\[
 \frac{P_*^2((u,v),(s,z))}{\Pi_*(s,z)}
 =\frac{T_*(v,s)}{\pi_*(s)}.
\]
For $(v,s)=(0,0),(0,1),(1,0),(1,1)$ these four ratios are
$3(1-\beta),\,3\beta/2,\,3q_*,\,3(1-q_*)/2$, respectively.
They are all at least $\beta+q_*\ge\beta$, checking every pair-state
certificate entry. The target at terminal $00$ is Bernoulli$(\beta)$,
not Bernoulli$(q_*)$.
For $2\le j\le T$ put $x^{(j)}=1^{M+2-j}0^j$.
These disjoint words have full stationary masses
\begin{equation}
 v_j=\frac23(1-q_*)^{M+1-j}q_*(1-\beta)^{j-1}
 \ge\frac1{8M}. \label{eq:null-word-masses}
\end{equation}
Bernoulli's inequality gives the prefix factor at least $1/2$ and
the suffix factor at least $7/8$, so in fact the product is at least
$7/(48M)>1/(8M)$. These masses include the initial symbol and every
one of the $M+1$ transitions.

\paragraph{A predeclared strict kernel with identity rejection.}
For $2\le j\le T$, define a finite measure on words by
\[
 \nu_j(dy)=\int_{At}^{1/4}
       Q_a(dy;F,U\ge u_0,J_{00}=j)\,\frac{da}{a},
 \qquad \mu_j=\nu_j(\text{all words}).
\]
Lemma~\ref{lem:mixture} gives $\mu_j\le t/u_0$.
For $\mu_j>0$, normalization $\nu_j/\mu_j$ is a probability law.
These laws depend only on declared $M,k,r$ and the auxiliary family,
not the true source, its parameters, its initial law, or a predictor seed.

On input $x^{(j)}$ with $\mu_j>0$, draw $Y$ once from
$\nu_j/\mu_j$. If $B$ holds, form the deterministic $C_r(Y)$ of
Lemma~\ref{lem:relocation}. Accept it only if its actual Hamming
distance from $x^{(j)}$ is at most $k$ and its final pair equals that
of $x^{(j)}$. Otherwise output the unchanged input. On zero-mass labels
and all other input words, also output the unchanged input. This defines
one source-independent kernel on all words, with a strict pointwise
budget and protected final pair, not merely an average budget.
It uses one draw, not repeated sampling until acceptance.

On $G$, Lemma~\ref{lem:relocation} bounds the distance by
$3k/16+4r\le3k/16+k/32<k$, so every check succeeds, and
\[
 \mathcal T_r(C_r(Y))=\mathcal T_r(Y).
\]
The same statistic identity holds on every transformed acceptance;
it is not asserted on identity rejections. These are the sampled
alternative's boundary windows, not necessarily the null input's
windows. Input-window edits are allowed and included in the Hamming
check. The suffix label obeys exactly
$J_{00}(x^{(j)})=j-2$ and $J_{00}(Y)=j$.

\paragraph{Nonnegative-loss ordering.}
For a sampled word $y$, let $\ell_a(y)$ be the predictor's expected
absolute error against $a$ on $\mathcal T_r(y)$, and let $\ell_*(y)$
be its corresponding error against $\beta$. On $F$ the clean actual
terminal target is $a$; the triangle inequality, also after averaging
independent decoder randomness, gives
$\ell_*(y)\ge a-\beta-\ell_a(y)$.
For $\mu_j>0$, \eqref{eq:null-word-masses} gives
\[
 \frac{v_j}{\mu_j}\ge\rho:=\frac{u_0}{8Mt}.
\]
First discard nonnegative losses from other null inputs and rejected
proposals. On the remaining nonnegative accepted losses, replace
$v_j/\mu_j$ by $\rho$. Next restrict to $G$, where acceptance and
statistic matching hold. Only then use the signed triangle bound:
\begin{align*}
 R_{\rm null}
 &\ge\rho\int_{At}^{1/4}
          \E_{Q_a}[\ind_{F\cap G}\ell_*(Y)]\,\frac{da}{a}\\
 &\ge\rho\int_{At}^{1/4}
    \left((a-\beta)Q_a(F\cap G)
          -\E_{Q_a}[\ind_{F\cap G}\ell_a(Y)]\right)\frac{da}{a}\\
 &\ge\rho\left(\frac t{16}-S\right)L_-.
\end{align*}
The restricted alternative loss is bounded by its unconditional clean
risk, hence by $S$. Risk is nonnegative, giving \eqref{eq:finite-lower}.
In particular, no lower mass multiplier has been placed in front of a
possibly negative quantity, and no endpoint-conditioned loss has been
substituted for the unconditional target risk.
\end{proof}

\subsection{Consequences for widths and supplemental messages}

\begin{proposition}[Width and fixed-fraction consequences]
\label{prop:lower-consequences}
For a prescribed sequence $2\le r_M\le M+2$, a finite uniform factor
relative to $g_M(k)=\max\{M^{-1/2},k/M\}$ exists if and only if
$\liminf_{M\to\infty}r_M/M>0$, for stationary or arbitrary starts.
For each fixed $0<\alpha\le1/4096$ and
$r_M=\max\{2,\lceil\alpha M\rceil\}$,
\[
 2^{-26}\log\frac1{32\alpha}
 \le C_{\rm stat}(\alpha)
 \le C_{\rm all}(\alpha)
 \le2^{19}\left(1+\log\frac2\alpha\right).
\]
With exact counts and protected terminal pair, optionally also the
initial pair, a post-corruption full-word encoder needs at least two
supplemental symbols to attain a finite uniform factor; two suffice
when its message is exact and fault-free.
\end{proposition}
\begin{proof}
Suppose first that the liminf is zero. Select an increasing subsequence
on which $r_M/M\to0$, and on that subsequence put
\[
 b_M=\max\{r_M,\sqrt M\},\qquad
 k_M=\lfloor\sqrt{Mb_M}\rfloor.
\]
Then $b_M=o(M)$, $k_M/M\to0$, $k_M/b_M\to\infty$,
$r_M/k_M\to0$, and $k_M/\sqrt M\to\infty$. Moreover
$M/k_M^2\to0$, so every condition in
\eqref{eq:finite-lower-domain} holds eventually. Here
$t_M=ck_M/M$, $\rho=u_0/(8Mt_M)$, $L_-\to\infty$, and
$\rho\to1/32$.
If one $k$-blind family had a finite admissible factor $C$, its clean
risk on every alternative at this same $M$ and supplied certificate
would satisfy $S_M\le C/\sqrt M=o(t_M)$. For example, eventually
$\rho\ge1/64$ and $t_M/16-S_M\ge t_M/32$, so
\[
 \frac{R_{\rm null}}{g_M(k_M)}
 \ge\frac{c}{2048}L_-\longrightarrow\infty,
\]
since $g_M(k_M)=k_M/M$ eventually. This contradicts $C$.
The alternatives may depend on the tested budget, but the predictor
and its clean guarantee are the same map at the same horizon. These
stationary witnesses also rule out an all-start guarantee.

Conversely, if the liminf is positive, choose $\delta>0$ and $M_0$
so $r_M\ge\delta M$ for $M\ge M_0$. Put
$m_M=\min\{r_M-2,\lfloor M/2\rfloor\}$.
Then $(m_M+1)/M$ is bounded away from zero for all sufficiently large
$M$. The suffix estimate of Proposition~\ref{prop:finite-upper} is at
most $2^{19}[1+\log(M/(m_M+1))]g_M(k)$ at $\beta=1/M$, so gives
one uniform all-start factor eventually. A finite initial set of
horizons is harmless: any prediction has risk at most one and
$g_M(k)\ge M^{-1/2}$. This proves the exact criterion without
monotonicity of the widths.

Now fix $\alpha>0$ in the stated interval before sending $M$ to
infinity. For any family with finite admissible stationary factor $C$,
put $k=128r_M$ at sufficiently large horizons. Then
\[
 \frac kM\to128\alpha\le\frac1{32},\quad
 t\to128c\alpha>0,\quad At\to8\alpha<\frac14,
\]
all finite-domain conditions eventually hold, and
\[
 L_-\to\log\frac1{32\alpha},\qquad
 \frac{u_0}{8Mt}\to\frac1{32},\qquad
 S/t\le\frac{C}{t\sqrt M}\to0.
\]
For these budgets $k\ge\sqrt M$ eventually. Using the deliberately
loose $g_M(k)\le2k/M$ in \eqref{eq:finite-lower}, the positive limiting
bracket gives
\[
 C\ge\liminf_{M\to\infty}\frac{R_{\rm null}}{g_M(k)}
 \ge\frac1{32}\frac{c}{32}\log\frac1{32\alpha}
 =2^{-26}\log\frac1{32\alpha}.
\]
This holds for every admissible finite factor, so taking the infimum
does not require an attained optimum. The upper follows from
Proposition~\ref{prop:finite-upper}: $m_M+1\ge\alpha M/2$, including
$r_M=2$, hence $1+\log(M/(m_M+1))\le1+\log(2/\alpha)$.
The stationary class is a subclass of the arbitrary-start class.
These arguments take the horizon limit for each fixed fraction and
factor; they neither exchange limits nor give an every-budget
logarithmic risk lower bound.

Finally, a one-symbol supplemental alphabet is constant, so leaves
only counts and the supplied endpoint pair or pairs. At $r_M=2$ even
both endpoint pairs are insufficient by the necessity just proved.
The exact one-bit factorization in
Proposition~\ref{prop:finite-upper} supplies two symbols and has
all-start risk at most $(2^{17}+2^{23})g_M(k)$. This concerns a
fault-free message computed by scanning the whole already corrupted
word, not one raw coordinate or a message subject to further corruption.
\end{proof}

\subsection{Elementary same-information oracle lower floors}

\begin{lemma}[Critical oracle lower floors]
\label{lem:oracle-floor}
For every $M\ge64$ and $0\le k\le M$, the stationary minimax risk
with supplied $\beta=1/M$ is at least $g_M(k)/96$, even if the predictor
knows $k$ and observes the entire corrupted word. Consequently this
lower bound holds for each restricted observation considered here and
for the arbitrary-start minimax risk.
\end{lemma}
\begin{proof}
We give separate clean and edit witnesses; taking their maximum is
permitted in a minimax lower bound.

\paragraph{Clean floor.}
Put $N=M+2$ and use stationary IID Bernoulli sources
\[
 p_\pm=\tfrac12\pm\frac1{16\sqrt N}.
\]
Their pair chains satisfy $P^2(u,v)=\Pi(v)$ for every row, hence the
supplied certificate. The target is $p_\pm$ at every terminal pair.
The elementary Bernoulli divergence bound
\[
 d(p\Vert q)\le\frac{(p-q)^2}{q(1-q)}
\]
follows by applying $\log x\le x-1$ to its two summands.
With $\Delta=p_+-p_-=1/(8\sqrt N)$, it gives one-letter divergence
at most $1/(8N)$. Independence makes the full-word divergence at most
$1/8$, and their total variation distance is at most $1/4$.
For clarity, this last inequality follows without a statistical
independence assertion about a selected row: for finite probability
laws $P,Q$, group their masses on
$A=\{x:P(x)\ge Q(x)\}$ and its complement. The log-sum inequality
(Jensen's inequality for $x\log x$) gives
$D(P\Vert Q)\ge d(P(A)\Vert Q(A))$.
For fixed $q$, the second derivative in $p$ of $d(p\Vert q)$ is
$1/[p(1-p)]\ge4$; its value and first derivative vanish at $p=q$.
Thus $D(P\Vert Q)\ge2(P(A)-Q(A))^2=2\TV(P,Q)^2$.

Testing by the nearer target makes each error incur prediction loss
at least $\Delta/2$. Every possibly randomized test has average error
at least $(1-\TV(P,Q))/2$, since the difference of its acceptance
probabilities is at most total variation. Thus the average prediction
risk is at least
\[
 \frac\Delta4(1-\TV(P,Q))
 \ge\frac3{128\sqrt N}\ge\frac1{64\sqrt M}.
\]
Identity corruption is allowed at every $k$, including for a
known-budget predictor.

\paragraph{Edit floor.}
For $k\ge1$, use the uniform prior on IID Bernoulli parameters
$p\in[1/4,3/4]$. Put $b=k+1$. If the first $M$ symbols contain $n$
ones, retain a uniformly chosen set of
$r=b\lfloor n/b\rfloor$ of those ones, turning the others to zero.
Leave the final two symbols $U$ unchanged. This single kernel is
source-independent and changes $n-r\le b-1=k$ editable bits on every
outcome.
For a specified retained set $T$ of size $r$, conditional on $n$,
\[
 \Pr(T\text{ retained}\mid n)
 =\frac{\binom{M-r}{n-r}}
        {\binom Mn\binom nr}
 =\frac1{\binom Mr}.
\]
Hence conditional on $r,U$ the output arrangement is uniform and
independent of $p$, even after mixing the possible $n$ in a bucket.
The full observation factors through at most
\[
 D=4\bigl(\lfloor M/b\rfloor+1\bigr)
\]
labels $(r,U)$ followed by a $p$-independent channel. The four protected
pair possibilities are counted, not conditioned away.

Realize that channel using a seed independent of $p$ and of the label,
and likewise use an independent decoder seed. Freeze both seeds.
There are at most $D$ possible prediction values. For the uniform
prior on an interval of length $\ell=1/2$, the union of radius-$v$
neighborhoods of those values has prior mass at most $2Dv/\ell$,
regardless of how a label is randomly selected conditional on $p$.
Outside that union every possible prediction has error greater than
$v$. Integrating this probability bound gives Bayes absolute-error
risk at least
\[
 \int_0^{\ell/(2D)}\left(1-\frac{2Dv}{\ell}\right)dv
 =\frac\ell{4D}.
\]
Averaging the independent seeds preserves it. Consequently some
parameter in the prior support has risk at least
\[
 \frac1{32(\lfloor M/(k+1)\rfloor+1)}
 \ge\frac{k+1}{32(M+k+1)}\ge\frac{k}{96M}.
\]
IID stationarity and $P^2(u,v)=\Pi(v)$ again give the same source
class and supplied certificate, with target $p$ at every row. Combining
the clean floor and edit floor yields $g_M(k)/96$, also at $k=0$.
A predictor with less information can only have a larger minimax risk,
and allowing arbitrary starts contains the stationary experiment.
\end{proof}

\section{Simultaneous finite-horizon competitive factors}
\label{sec:finite-window-proof}

We prove Theorem~\ref{thm:finite-window} using the finite estimates proved
in the preceding two sections. No limit in $M$ or $r$ is taken. Throughout
this section $\ell=1+\log((M+2)/r)\ge1$, the certificate is $1/M$, and
$\Gamma$ denotes the finite-horizon quantities of
\eqref{eq:finite-factor}, not the sequence factors $C$.

\begin{proof}[Proof of Theorem~\ref{thm:finite-window}]
\noindent\textbf{Upper bound, including overlapping windows.}
Set $m=\min\{r-2,\lfloor M/2\rfloor\}$ and
$L=1+\log(M/(m+1))$. Since $r-1\ge r/2$ and
$\lfloor M/2\rfloor+1\ge M/2$,
\[
 \frac{M}{m+1}
 \le\max\left\{\frac{2M}{r},2\right\}
 \le\frac{2(M+2)}r.
\]
The last step uses $r\le M+2$ and covers full or overlapping windows.
Hence $L\le\ell+\log2\le2\ell$. Proposition~\ref{prop:finite-upper}
gives one counts-and-suffix map, independent of $k$, with all-start risk
at most $2^{19}Lg_M(k)\le2^{20}\ell g_M(k)$ for every $k$.
Stationary experiments are a subclass of arbitrary-start experiments,
so taking the infima over the same maps also gives
$\Gamma_{\rm stat}(M,r)\le\Gamma_{\rm all}(M,r)$.

\noindent\textbf{Lower bound and the small-logarithm case.}
Fix any allowable map $A$ and write
$G=\Gamma_{\rm stat}(A;M,r)\le\sqrt M$.
Lemma~\ref{lem:oracle-floor}, at $k=0$, implies $G\ge1/96$ even if
full history were available. If $\ell<64$, then
$1/96\ge2^{-30}\ell$, proving the desired bound. For $\ell\ge64$,
there is nothing to prove when $G\ge2^{-30}\ell$. Suppose instead
\begin{equation}\label{eq:fh-contradiction}
 G<2^{-30}\ell.
\end{equation}
As $r\ge2$ and $M\ge2$, $(M+2)/r\le M$. Thus
$\log M\ge\ell-1\ge63$, and $M>2^{80}$ because $80\log2<56<63$.
Choose the integer budget
\begin{equation}\label{eq:fh-budget}
 k=\left\lceil\max\{128r,M^{3/4}\}\right\rceil.
\end{equation}
This budget depends only on $(M,r)$; the same map $A$ is used at this
budget and at zero.

\noindent\textbf{Every finite lower-domain condition.}
Use the constants from Proposition~\ref{prop:finite-lower}, writing
$A_0=2^{12}$ for its numerical constant to distinguish it from the map:
\[
 c=2^{-16},\qquad t=ck/M,\qquad
 u_0=\lfloor ck/4\rfloor,\qquad
 L_-=\log\frac1{4A_0t}.
\]
The choice of $k$ gives $r\le k/128$, $k\ge M^{3/4}$, and
\[
 \frac M{k^2}\le M^{-1/2}<2^{-40}
   =\frac{A_0c^2}{2^{20}},\qquad
 ck\ge2^{-16}M^{3/4}>2^{44}\ge1024.
\]
Moreover
\[
 \frac rM=(1+2/M)e^{1-\ell}\le2e^{-63}<2^{-79}.
\]
Since $k\le128r+M^{3/4}+1$, this implies
\[
 \frac kM<2^{-72}+2^{-20}+2^{-80}<2^{-19}<\frac18.
\]
In particular $k$ is a legal budget, $k\le M/8$, and
$A_0t=k/(16M)<1/128<1/4$. Together with the original integer
$M\ge64,r\ge2$, these verify every condition in
\eqref{eq:finite-lower-domain}. The windows needed in that proof are
nonoverlapping since $r\le k/128\le M/1024$.
Also $ck\ge1024$ implies $u_0\ge ck/8$, so
\begin{equation}\label{eq:fh-prefactor}
 \rho:=\frac{u_0}{8Mt}\ge\frac1{64}.
\end{equation}

\noindent\textbf{The same map's clean risk.}
Let $S$ be the supremum of the clean risks on the stationary alternatives
in Proposition~\ref{prop:finite-lower}. All qualify at the same supplied
certificate, so $S\le G/\sqrt M$. For $x\ge0$, differentiation shows
that $(1+x)e^{-x/4}$ is maximized at $x=3$, with value
$4e^{-3/4}<4$. Thus $\ell\le1+\log M\le4M^{1/4}$ and
\eqref{eq:fh-contradiction} gives the finite inequalities
\begin{equation}\label{eq:fh-clean}
 S<2^{-28}M^{-1/4}\le2^{-21}M^{-1/4}\le t/32.
\end{equation}
The last step uses $k\ge M^{3/4}$ and $c/32=2^{-21}$.
No asymptotic small-order estimate or distinct clean estimator is used.

\noindent\textbf{A logarithm uniform in both parameters.}
Since $4A_0c=1/4$, $L_-=\log(4M/k)$. Put
$x=\max\{128r,M^{3/4}\}\ge1$; then $k=\lceil x\rceil\le2x$, and
\[
 L_-\ge\min\left\{\log\frac M{64r},\ \log(2M^{1/4})\right\}.
\]
The first term equals
$\ell-1-\log(1+2/M)-\log64$ and is at least
$\ell-2-\log64\ge\ell-8\ge\ell/8$, using
$\log(1+2/M)\le1$, $\log64<6$, and $\ell\ge64$.
The second term is at least $(\ell-1)/4\ge\ell/8$.
Consequently $L_-\ge\ell/8$.

\noindent\textbf{Common-kernel comparison and infimum.}
Proposition~\ref{prop:finite-lower} supplies a qualifying stationary null
and one strict source-independent kernel, preserving the actual final pair,
with
\[
 R_{\rm null}\ge\rho\left[\left(t/16-S\right)L_-\right]_+.
\]
Its full proof above includes source legality, mixture normalization,
window matching, and pointwise budget checks. Here the bracket is positive
by \eqref{eq:fh-clean}. Since $k\ge M^{3/4}\ge\sqrt M$, one has
$g_M(k)=k/M$. Therefore
\[
 G\ge\frac{R_{\rm null}}{g_M(k)}
 \ge\frac{(1/64)(t/32)(\ell/8)}{k/M}
 =2^{-30}\ell,
\]
contradicting \eqref{eq:fh-contradiction}. The lower bound thus holds for
every map and passes to its infimum without requiring attainment.
The witnesses are stationary and hence are also available in the all-start
class. This completes the finite-horizon proof.
\end{proof}

For clarity about the sequence consequences, if $r_M/M$ has positive
liminf, then $\ell_M$ is bounded eventually, and finitely many earlier
horizons have finite factor; the explicit suffix maps give a single
uniform family. If that liminf is zero, $\ell_M$ diverges on a subsequence,
so the finite lower forbids a uniform factor for any family. For fixed
$\alpha$, the finite bounds give the logarithmic order, while the original
fixed-$\alpha$ proof in Proposition~\ref{prop:lower-consequences} retains
its stronger $2^{-26}\log(1/(32\alpha))$ lower constant. Neither argument
claims a logarithmic risk lower at every individual budget.

\section{A direct-count approximation interface}
\label{sec:approximate-proofs}

This section changes the representation of the half-gate statistics, not
its statistical mechanism.  An earlier twelve-coordinate construction
already established the real-error half-gate argument and the finite-bank
precision, noise, and source-edit phenomenon \citep{BCJ001}.  Here the supplied
coordinates directly approximate $J,C,D$, rather than reconstructing $J$
from two nonnegative successor coordinates.  Sequential clipping makes
this changed interface legal.  In particular, the result is not a claim
that exact recovery of a discontinuous gate bit is necessary, nor a new
precision principle.

\subsection{Statement and clipping}

For every row $u=(a,b)$, retain $d(u)=1-a$, and define the clean counts
\begin{align}
 J_u&=\sum_{i=0}^{M-1}\ind\{Y_i=u\},&
 C_u&=\sum_{i=0}^{M-1}\ind\{Y_i=u,\ X_{i+2}=d(u)\},
 \label{eq:a-full-counts}\\
 D_u&=\sum_{i=0}^{\lfloor M/2\rfloor-1}\ind\{Y_i=u\}.
 \label{eq:a-half-count}
\end{align}
Thus $D$ denotes a first-half departure count, always normalized by $M$
when used as a frequency.  Primed counts are evaluated on the corrupted
word $Z$.  At the actual protected row $u=Y_M$, omit the subscript and
write $q=\theta_u(d(u))$.  The protected pair makes the selected row the
same for the clean and corrupted words.

For real $x$ and $a\le b$, write
$\operatorname{clip}(x;a,b)=\min\{b,\max\{a,x\}\}$.
Suppose an observation-measurable computation supplies, for every row,
real values $\widetilde J_u,\widetilde C_u,\widetilde D_u$.  Its only
accuracy premise is the simultaneous pathwise bound
\begin{equation}\label{eq:a-stat-error}
 \max_u\max\{
 |\widetilde J_u-J'_u|,
 |\widetilde C_u-C'_u|,
 |\widetilde D_u-D'_u|\}\le e,
 \qquad e\ge0.
\end{equation}
This must hold on every allowed word and for every outcome of the
computation's independent random coins.  The computation has no access
to the clean word or hidden source parameters.  Its errors may be
dependent across coordinates and may depend on the observed word.
The number $e$ is a proved error bound, not an input required by the
decoder.  Clip sequentially, using the already clipped first coordinate:
\begin{equation}\label{eq:a-sequential-clip}
 \overline J_u=\operatorname{clip}(\widetilde J_u;0,M),\qquad
 \overline C_u=\operatorname{clip}(\widetilde C_u;0,\overline J_u),\qquad
 \overline D_u=\operatorname{clip}(\widetilde D_u;0,\overline J_u).
\end{equation}
The approximate half-gate returns designated probability
\begin{equation}\label{eq:a-gate}
 \widehat q=
 \begin{cases}
 \overline C/\overline J,
       &\overline J>0\ \text{and}\ 512\overline D\ge\overline J,\\
 0,&\text{otherwise}.
 \end{cases}
\end{equation}
Equality accepts.  A final observation-measurable probability
$\widehat q_{\rm out}\in[0,1]$ with
$|\widehat q_{\rm out}-\widehat q|\le\xi$, $\xi\ge0$, is also allowed.
The complementary symbol receives probability $1-\widehat q_{\rm out}$.

\begin{theorem}[Approximate direct half-gate statistics]
\label{thm:approximate-gate}
For every integer $M\ge64$, every integer $0\le k\le M$, every homogeneous
binary order-two source with an all-state certificate
$P^2(u,v)\ge\beta\Pi(v)$, $\beta>0$, every initial law, and every legal
source-independent edit kernel in \eqref{eq:attack}, the decoder
\eqref{eq:a-sequential-clip}--\eqref{eq:a-gate}, with the pathwise premise
\eqref{eq:a-stat-error}, has unsquared total-variation risk
\begin{equation}\label{eq:a-risk-general}
 R\le\min\left\{1,
 \frac{65536}{\sqrt M}+\frac{65536}{M\sqrt\beta}
 +1072176\frac{6k+e}{M}+\xi\right\}.
\end{equation}
The same decoder works simultaneously for all these parameters without
knowing $k,\beta$, or $e$.  At the supplied certificate $\beta=1/M$,
\begin{equation}\label{eq:a-risk-critical}
 R\le\min\left\{1,
 \frac{131072}{\sqrt M}+6433056\frac{k}{M}
 +1072176\frac eM+\xi\right\}.
\end{equation}
The edit kernel does not observe the predictor's random seed, and the
result permits arbitrary independent randomization satisfying the
pathwise accuracy premise.
\end{theorem}

\begin{lemma}[Clipping with a perturbed upper endpoint]
\label{lem:a-clipping}
Let $0\le c\le j\le M$ and
$|\widetilde j-j|,|\widetilde c-c|\le e$.
Set $\overline j=\operatorname{clip}(\widetilde j;0,M)$ and
$\overline c=\operatorname{clip}(\widetilde c;0,\overline j)$.
Then $|\overline j-j|\le e$ and $|\overline c-c|\le e$.
The same statement holds with $c$ replaced by any $d\in[0,j]$.
\end{lemma}
\begin{proof}
Projection onto $[0,M]$ does not increase distance to $j$.
For the second coordinate,
$\overline c\le\max\{\widetilde c,0\}\le c+e$.
For its lower bound, if $\widetilde c<0$, then $c\le e$, so
$\overline c=0\ge c-e$.  If
$0\le\widetilde c\le\overline j$, clipping retains
$\widetilde c\ge c-e$.  If $\widetilde c>\overline j$, then
$\overline c=\overline j\ge j-e\ge c-e$.
These cases exhaust the possibilities, including $\overline j=0$;
no clipped coordinate is required to be integer.
\end{proof}

\subsection{Complete pathwise perturbation argument}

One symbol edit affects at most three departure triples.  Changing a
triple has histogram $\ell_1$ cost at most two, so each full-count
coordinate changes by at most $6k$.  Restricting to the first half does
not increase this bound.  Adjacent edits and triples crossing the split
are covered by the union bound; their affected triples are not assumed
disjoint.  Lemma~\ref{lem:a-clipping} therefore gives, on every edit and
computation outcome,
\begin{equation}\label{eq:a-clean-errors}
 |\overline J-J|,
 |\overline C-C|,
 |\overline D-D|\le\Delta:=6k+e.
\end{equation}
The perturbation radius $\Delta$ may be an arbitrarily small noninteger.
Define the events
\[
 G=\{\overline J>0,\ 512\overline D\ge\overline J\},\qquad
 H=\{D\ge J/1024\},\qquad
 B=\{J>0,\ D<J/256\}.
\]
Both $C/J$, where defined, and the accepted clipped ratio belong to
$[0,1]$.

On $G\cap H\cap\{J>0\}$, we claim
\begin{equation}\label{eq:a-ratio}
 \left|\frac{\overline C}{\overline J}-q\right|
 \le\left|\frac CJ-q\right|+\frac{8\Delta}{J+1}.
\end{equation}
If $\Delta=0$, the ratios coincide.  If $\Delta>0$ and $J\ge2\Delta$,
then $\overline J\ge J/2$, and subtracting ratios gives
\[
 \left|\frac{\overline C}{\overline J}-\frac CJ\right|
 \le\frac{|\overline C-C|}{\overline J}
       +\frac{C|J-\overline J|}{J\overline J}
 \le\frac{4\Delta}{J}
 \le\frac{8\Delta}{J+1},
\]
where the last step uses the integer inequality $J\ge1$.
If $0<J<2\Delta$, then $\Delta>J/2$ and
$8\Delta/(J+1)\ge1$, so the unit bound on the ratio difference suffices.
This proves \eqref{eq:a-ratio} for all real $\Delta\ge0$.

On $G\cap H^c$, acceptance and \eqref{eq:a-clean-errors} imply
\[
 D+\Delta\ge\overline D
 \ge\frac{\overline J}{512}\ge\frac{J-\Delta}{512},
 \qquad D<\frac J{1024}.
\]
Combining them yields $J<1026\Delta$.
On $G^c\cap B^c\cap\{J>0\}$, if $\overline J=0$ then
$J\le\Delta$.  Otherwise the strict rejection inequality and
$D\ge J/256$ give
\[
 \frac J{256}-\Delta
 \le\overline D<\frac{\overline J}{512}
 \le\frac{J+\Delta}{512},
\]
and hence $J<513\Delta$.  For $\Delta=0$, neither false-status case
with $J>0$ is possible.  Rejection on $B$ has loss exactly $q$.

The occupation lemma has an integer threshold, so introduce
\begin{equation}\label{eq:a-ceilings}
 w_1=\max\{1,\lceil1026\Delta\rceil\},\qquad
 w_2=\max\{1,\lceil513\Delta\rceil\}.
\end{equation}
They cover the strict false-status inequalities and $J=0$.
In the separate $\overline J=0$, $J>0$ case, $J\le\Delta$ forces
$\Delta>0$, and then $J<513\Delta$ as required.
For every $\Delta\ge0$,
$w_1+w_2\le1539\Delta+2$, including $\Delta=0$.
The preceding cases, bounding any remaining loss by one, prove the
pointwise reduction
\begin{equation}\label{eq:a-pathwise}
 \begin{split}
 |\widehat q-q|\le{}&
 \ind\{J>0\}\left|\frac CJ-q\right|
 +8\Delta\frac{\ind_H}{J+1}+q\ind_B\\
 &+\ind\{J<w_1\}+\ind\{J<w_2\}.
 \end{split}
\end{equation}
The ratio $C/J$ is evaluated only when $J>0$.  At $J=0$ the tail
indicators cover the loss without defining that ratio.  All right-hand
terms are nonnegative.  In particular, no consistency assumption that
the approximate coordinates are counts of a word is used, and no exact
agreement with the corrupted word's gate bit is asserted.

\subsection{Averaging and normalized precision}

\begin{proof}[Proof of Theorem~\ref{thm:approximate-gate}]
The four stochastic inputs have complete proofs in
Section~\ref{sec:positive-proofs}.  With $D=J_{1,Y_M}$, they are precisely
\begin{align*}
 \Prb(J<w)&\le\min\{1,16w/(M+1)\}
          &&\text{for integer }w\ge1
          &&\text{(Lemma~\ref{lem:occupation})},\\
 \E\!\left[\ind\{J>0\}\left|C/J-q\right|\right]
      &\le512/\sqrt{M+1}
          &&&&\text{(Lemma~\ref{lem:empirical})},\\
 \E[q\ind_B]&\le65536/(M\sqrt\beta)
          &&&&\text{(Lemma~\ref{lem:imbalance})},\\
 \E[\ind_H/(J+1)]&\le130944/M
          &&&&\text{(\eqref{eq:p-half-inverse} in
                    Proposition~\ref{prop:finite-upper})}.
\end{align*}
In particular these inputs hold for arbitrary starts, including starts
outside stationary support.  They require neither endpoint-conditioned
independent departures nor independence of the perturbation errors.
Apply them to \eqref{eq:a-pathwise}, then use \eqref{eq:a-ceilings}:
\[
 \begin{split}
 \E|\widehat q-q|\le{}&
 \frac{512}{\sqrt{M+1}}+\frac{65536}{M\sqrt\beta}
 +\frac{8\cdot130944\Delta}{M}
 +\frac{16(1539\Delta+2)}{M+1}\\
 \le{}&\frac{65536}{\sqrt M}+\frac{65536}{M\sqrt\beta}
       +\frac{1072176\Delta}{M}.
 \end{split}
\]
Here
\[
 8\cdot130944+16(1026+513)=1072176,\qquad
 \frac{512}{\sqrt{M+1}}+\frac{32}{M+1}
 \le\frac{65536}{\sqrt M}.
\]
The final output error contributes at most $\xi$ by the triangle
inequality.  A legal binary probability has loss at most one.
Substitute $\Delta=6k+e$ to prove \eqref{eq:a-risk-general};
then $\beta=1/M$ proves \eqref{eq:a-risk-critical}.
The proof is pointwise before averaging over all permissible
randomizations, which proves the stated simultaneous quantifiers.
\end{proof}

For a nonnegative integer $h$, let $Q=2^h$ and separately round each
of the twelve normalized primed statistics
$J'_u/M,C'_u/M,D'_u/M$ to its nearest point of
$\{0,1/Q,\ldots,1\}$, with any fixed tie rule.  Multiplication by $M$
gives approximations satisfying $e\le M/(2Q)$.  Thus
\eqref{eq:a-risk-critical} gives the explicit rounded-bank upper bound
\begin{equation}\label{eq:a-rounded-risk}
 R\le\min\left\{1,
 \frac{131072}{\sqrt M}+6433056\frac kM+\frac{536088}{Q}+\xi
 \right\}.
\end{equation}
Consequently $h\ge\lceil\tfrac12\log_2 M\rceil$ and
$\xi\le M^{-1/2}$ suffice for a constant multiple of $g_M(k)$,
uniformly in the unknown budget.  If a separate output stage has error
at most $2^{-b}$, its contribution is simply $\xi=2^{-b}$; the output
precision $b$ is distinct from the bank precision $h$ and the raw-window
width $r$.

For known $M$, a straightforward integer-numerator representation of
this bank uses at most $12(h+1)$ bits, plus the two exact protected
endpoint bits.  The extra numerator bit permits the grid value one.
This is a representation statement, not a streaming algorithm with
that workspace: rounding after each update can accumulate error.
Nor does it prove that all twelve coordinates are necessary.
Raw input storage, exact prefix selection, addressing, accumulation,
pair/triple feature construction, division, output precision, weights,
depth, attention work, and training are separate costs.  In particular,
a classical counter has the same statistical interface; there is no
neural statistical advantage in \eqref{eq:a-rounded-risk}.

\section{Precision of the fixed nearest-rounded positional bank}
\label{sec:precision-proofs}

The lower bound in this section concerns a particular observation, not
an arbitrary encoder with the same number of bits.  The earlier
finite-bank result already established a twelve-coordinate precision
floor together with noise and source-edit bounds \citep{BCJ001}.  Its floored
triple/half-count coordinates differ from the nearest-rounded direct
$J/C/D$ coordinates below.  Before quantization the two representations
are linearly interconvertible, but separate coordinatewise quantization
does not preserve that equivalence.  We therefore give a complete
lower proof for this changed cut.  Its common-cell and two-point
argument also includes $h=0$; it does not assert a new
statistical/quantization transition.

\subsection{The exact observation and result}

Fix integers $M\ge64$ and $h\ge0$, and set $Q=2^h$.
Let $\mathcal N_Q$ round to the nearest point of
$\{0,1/Q,\ldots,1\}$ using a fixed tie rule.
Use the counts in \eqref{eq:a-full-counts}--\eqref{eq:a-half-count},
and form the twelve-coordinate bank
\begin{equation}\label{eq:p-bank}
 \mathcal B_{M,h}(X)=
 \bigl(\mathcal N_Q(J_u/M),\mathcal N_Q(C_u/M),
       \mathcal N_Q(D_u/M)\bigr)_{u\in\{00,01,10,11\}}.
\end{equation}
The decoder observes exactly
\begin{equation}\label{eq:p-bank-cut}
 \mathcal O=(\mathcal B_{M,h}(X),E),\qquad
 E=(X_M,X_{M+1}),
\end{equation}
and source-independent public format metadata.  The endpoint is exact.
No raw history, scratch counts, other token banks, rounding residuals,
or additional feature port crosses this cut.

Let $\mathcal C_M$ be the clean class of homogeneous binary order-two
sources with arbitrary initial law admitting a stationary pair law
$\Pi$ with $P^2(u,v)\ge\Pi(v)/M$ for every $u,v$.  Define
\begin{equation}\label{eq:p-bank-minimax}
 R^*(M,h)=\inf_A\sup_{(\theta,\mu)\in\mathcal C_M}
 \E\,\TV\bigl(A(\mathcal O),\theta_{Y_M}\bigr).
\end{equation}
The infimum permits arbitrary computation, any binary probability
output, and arbitrary independent decoder seeds.  It is not restricted
to the gate decoder or to a finite output grid.  Public randomization
independent of the source and sample may be included in these seeds;
source-dependent advice or randomization correlated with hidden history
is not part of \eqref{eq:p-bank-cut}.

\begin{theorem}[Fixed-bank precision]
\label{thm:bank-precision}
For every integer $M\ge64$ and every integer $h\ge0$,
\begin{equation}\label{eq:p-bank-rate}
 \frac1{1600000}\max\{M^{-1/2},2^{-h}\}
 \le R^*(M,h)
 \le\min\bigl\{1,667160\max\{M^{-1/2},2^{-h}\}\bigr\}.
\end{equation}
Both lower components already hold on the stationary IID Bernoulli
subclass with parameter $p\in[1/8,1/4]$.  Consequently the same lower
bounds hold for an adversarial experiment that allows zero edits,
provided its decoder has no more information than the fixed cut
\eqref{eq:p-bank-cut}.
\end{theorem}

\subsection{Two elementary probability inequalities}

For completeness, we give the general concentration and testing
inequalities needed below.  They apply to finite input laws; randomized
decoder seeds will be handled separately without restricting their
support.

\begin{lemma}[Bounded differences]
\label{lem:p-bounded-differences}
Let $F$ be a real function of independent finite-valued coordinates,
and suppose changing coordinate $i$ alone changes $F$ by at most
$c_i\ge0$.  With $V=\sum_i c_i^2>0$, for every $t>0$,
\[
 \Prb(|F-\E F|\ge t)\le2\exp(-2t^2/V).
\]
If $V=0$, $F$ is constant on the input support.
\end{lemma}
\begin{proof}
Reveal the independent coordinates in order and let $Z_i$ be the
increments of the Doob martingale of $F$.  Conditional on any revealed
prefix, $Z_i$ has mean zero and lies in an interval of length at most
$c_i$.  Indeed, compare two values of coordinate $i$ by integrating
over the same independent remaining coordinates; the resulting
conditional expectations differ by at most $c_i$.

For a centered random variable $Z$ in an interval $[a,b]$ of length
$c$, put $g(\lambda)=\log\E e^{\lambda Z}$.  Its second derivative
is the variance under the exponentially tilted law.  For any law on
$[a,b]$, the pointwise inequality $(Z-a)(b-Z)\ge0$ implies
\[
 \Var(Z)\le(\E Z-a)(b-\E Z)\le(b-a)^2/4.
\]
The same bound therefore holds under every tilt.  Since
$g(0)=g'(0)=0$, integrating twice gives
$g(\lambda)\le\lambda^2c^2/8$ for either sign of $\lambda$.
Apply this bound conditionally to each $Z_i$ and iterate conditional
expectations to obtain
\[
 \E e^{\lambda(F-\E F)}\le
 \exp\!\left(\lambda^2\sum_i c_i^2/8\right).
\]
Markov's inequality, optimized at $\lambda=4t/V$, gives
$\Prb(F-\E F\ge t)\le e^{-2t^2/V}$.
Apply the same argument to $-F$ and add the two tails.
The zero-sensitivity case is immediate.
\end{proof}

\begin{lemma}[Pinsker's inequality for finite laws]
\label{lem:p-pinsker}
For finite probability laws $P,Q$,
\[
 \TV(P,Q)\le\sqrt{\KL(P\Vert Q)/2}.
\]
Moreover, relative entropy adds across independent product factors.
\end{lemma}
\begin{proof}
If $P$ has positive mass where $Q$ has zero mass, the divergence is
infinite and the bound is immediate.  Otherwise set
$A=\{x:P(x)>Q(x)\}$, so
$\TV(P,Q)=P(A)-Q(A)$.
For a subset with $P$-mass $a>0$ and $Q$-mass $b$, its contribution to
relative entropy is $a$ times an expectation of
$\log(P(x)/Q(x))$ under normalized $P$ weights.  Jensen's inequality
for the convex function $-\log$ bounds this below by
$a\log(a/b)$: indeed the normalized expectation of $Q(x)/P(x)$
is at most $b/a$, with equality if $Q$ places no mass on zeros of $P$
inside the subset.  Zero masses follow by continuity or infinite
divergence.  Apply this log-sum inequality to $A$ and its complement:
\[
 \KL(P\Vert Q)\ge\operatorname{kl}(P(A)\Vert Q(A)).
\]
For fixed $q\in(0,1)$, the binary divergence
$f(p)=\operatorname{kl}(p\Vert q)$ satisfies
$f(q)=f'(q)=0$ and $f''(p)=1/[p(1-p)]\ge4$.
Integrating twice from $q$ to $p$, in either direction, gives
$f(p)\ge2(p-q)^2$.  Boundary cases follow by continuity or infinite
divergence.  Thus $\KL(P\Vert Q)\ge2\TV(P,Q)^2$.
Finally, a product-density ratio is a product of the factor ratios;
expanding its logarithm into a sum and taking expectations proves
additivity for independent products.
\end{proof}

\subsection{A common rounding cell for the complete mean vector}

Consider stationary IID Bernoulli$(p)$ input symbols, with
$p\in I=[1/8,1/4]$, and put $t=\lfloor M/2\rfloor/M$.
The means of the unrounded normalized coordinates, in row order
$00,01,10,11$, are
\begin{align}
 \E_p(J_u/M)_u
   &=((1-p)^2,p(1-p),p(1-p),p^2),\label{eq:p-bank-J-means}\\
 \E_p(C_u/M)_u
   &=(p(1-p)^2,p^2(1-p),p(1-p)^2,p^2(1-p)),
     \label{eq:p-bank-C-means}\\
 \E_p(D_u/M)_u
   &=t((1-p)^2,p(1-p),p(1-p),p^2).
     \label{eq:p-bank-D-means}
\end{align}
Every summand has the displayed law, including the final triple ending
at $X_{M+1}$.  Overlapping pairs and triples need not be independent
for these identities, which use only linearity of expectation.
The half bank has $\lfloor M/2\rfloor$ summands and denominator $M$.
For $M\ge64$, $t\ge1/3$.

On $I$, the four $J$-coordinate derivative magnitudes are bounded below
by $3/2,1/2,1/2,1/4$, respectively.  The derivatives of the two
distinct $C$ polynomials are
\[
 (1-p)(1-3p)\ge3/16,\qquad
 p(2-3p)\ge13/64.
\]
The $D$ derivatives are $t$ times the corresponding $J$ derivatives.
Thus all twelve mean functions are strictly monotone on $I$, and every
derivative has magnitude at least $d=1/12$.  Duplicated mean functions
do not cause a problem, and none of the coordinates is constant.

The rounding thresholds are
\[
 \mathcal T_Q=\{(j+1/2)/Q:0\le j\le Q-1\}.
\]
Each monotone coordinate mean crosses at most $Q$ thresholds.
Take the distinct crossing points in the interior of $I$ and include
the two endpoints.  These partition $I$ into at most $12Q+1$ open
intervals.  One such interval $(a_*,b_*)$ has length
\begin{equation}\label{eq:p-cell-length}
 L=b_*-a_*\ge\frac{1/8}{12Q+1}\ge\frac1{104Q}.
\end{equation}
Choose $p_0=a_*+L/3$ and $p_1=a_*+2L/3$.  Then
\begin{equation}\label{eq:p-cell-separation}
 \delta_p:=p_1-p_0=L/3\ge\frac1{312Q},\qquad
 \delta_p\le\frac1{24}.
\end{equation}
The mean vectors at both parameters round to one common bank
$\mathbf b_*$.  More precisely, for each coordinate mean $f$ and
$j\in\{0,1\}$,
\begin{equation}\label{eq:p-cell-margin}
 \operatorname{dist}(f(p_j),\mathcal T_Q)
 \ge dL/3\ge\frac1{3744Q}=:\rho.
\end{equation}
To justify this at the first and last partition intervals as well,
no threshold lies strictly between the values $f(a_*)$ and $f(b_*)$.
A threshold outside this value interval is separated from $f(p_j)$
by at least its distance to the corresponding endpoint value; integrating
the derivative lower bound over the relevant subinterval gives at
least $dL/3$.  This also handles thresholds outside $f(I)$, without
extending the derivative bound beyond $I$.  Endpoint ties are harmless
because $p_0,p_1$ are interior.  The argument includes $h=0$, when
the sole threshold is $1/2$.

\subsection{Concentration of the actual bank and its joint endpoint law}

Let $F$ be any one of the twelve unrounded normalized coordinates.
Changing one raw bit affects at most three relevant departure triples,
so changes the count by at most three.  Pair and half-pair coordinates
in fact require at most two.  There are $N=M+2$ independent raw bits,
so the bounded-difference constants satisfy
\[
 \sum_i c_i^2\le(M+2)(3/M)^2\le10/M,\qquad M\ge64.
\]
Lemma~\ref{lem:p-bounded-differences} yields
\[
 \Prb_p(|F-\E_pF|\ge\rho)\le2\exp(-M\rho^2/5).
\]
If every coordinate is strictly within $\rho$ of its mean, the actual rounded
bank equals $\mathbf b_*$.  The union bound and
\eqref{eq:p-cell-margin} give, separately at $p_0$ and $p_1$,
\begin{equation}\label{eq:p-bank-concentration}
 \Prb_{p_j}(\mathcal B_{M,h}\ne\mathbf b_*)
 \le\epsilon:=24\exp\!\left(-\frac{M}{5\cdot3744^2Q^2}\right).
\end{equation}
This is concentration of the random bank, not a substitution of the
rounded expectation for the expectation of a rounded variable.
No independence between bank coordinates or count summands is needed.

Set $C_0=500000000$.  If $M\ge C_0Q^2$, then
$\epsilon\le1/16$.  For an elementary numerical verification,
$C_0>30\cdot3744^2$, and the nonnegative Taylor terms for $e^6$
through degree ten already exceed $384$.  Hence
$24e^{-C_0/(5\cdot3744^2)}<24/384=1/16$.

Let $P_j$ be the joint law of $(\mathcal B_{M,h},E)$ at $p_j$, and
let $\widetilde P_j$ be the law of $(\mathbf b_*,E)$ with the same
endpoint marginal.  Couple them by retaining exactly the same endpoint
$E$ and replacing only the bank.  Equation~\eqref{eq:p-bank-concentration}
implies
\[
 \TV(P_j,\widetilde P_j)\le\epsilon.
\]
Under IID sampling, the endpoint marginal is
$\operatorname{Bern}(p_j)^{\otimes2}$.  Couple the two endpoint bits
using common independent uniforms at the two parameters.  A union
bound gives $\TV(\widetilde P_0,\widetilde P_1)\le2\delta_p$.
Therefore, when $M\ge C_0Q^2$,
\begin{equation}\label{eq:p-bank-TV}
 \TV(P_0,P_1)\le2\epsilon+2\delta_p
 \le\frac18+\frac1{12}=\frac5{24}<\frac12.
\end{equation}
The bank is not asserted independent of its endpoint: the final
departures explicitly use the protected bits.  Keeping the same
endpoint in the first coupling is what validates this argument.

\subsection{Randomized two-point loss and the full-history floor}

For any decoder let $A$ denote its probability assigned to symbol one.
In the IID submodel the true probability is $p$ at every actual
endpoint, so total variation loss is exactly $|A-p|$.
For two finite observation laws $P_0,P_1$, their common
subprobability measure has mass $1-\TV(P_0,P_1)$.
Append the same independent seed law to both observations, allowing
arbitrary seed spaces.  This preserves total variation and gives a
common subprobability product measure of the same mass.  On every
observation and seed,
$|A-p_0|+|A-p_1|\ge|p_1-p_0|$.
Integrating this inequality against that common measure, and using
nonnegativity to bound the two risks below, proves
\begin{equation}\label{eq:p-two-point}
 \max_{j=0,1}\E_j|A-p_j|
 \ge\frac{|p_1-p_0|}{2}\bigl(1-\TV(P_0,P_1)\bigr).
\end{equation}
This covers arbitrary independent randomized decoders and public
source-independent randomization.  Apply \eqref{eq:p-bank-TV} and
\eqref{eq:p-cell-separation} to obtain
\begin{equation}\label{eq:p-large-regime}
 R^*(M,h)\ge\frac{\delta_p}{4}\ge\frac1{1248Q},
 \qquad M\ge C_0Q^2.
\end{equation}

For the complementary regime, we prove a floor even when the decoder
sees all $N=M+2$ raw bits.  Take
$p_0=1/8$ and $p_1=p_0+\delta$ with
$\delta=1/(16\sqrt N)$, so both belong to $I$.
The inequality $\log x\le x-1$ gives the Bernoulli divergence bound
\[
 \begin{split}
 \KL(\operatorname{Bern}(p_0)\Vert\operatorname{Bern}(p_1))
 &\le p_0\left(\frac{p_0}{p_1}-1\right)
 +(1-p_0)\left(\frac{1-p_0}{1-p_1}-1\right)\\
 &=\frac{\delta^2}{p_1(1-p_1)}
 \le\frac{64}{7}\delta^2.
 \end{split}
\]
The last bound uses $p_1\in I$, where $p_1(1-p_1)\ge7/64$.
Product additivity in Lemma~\ref{lem:p-pinsker} bounds the full-history
divergence by $N(64/7)\delta^2=1/28$, and therefore its total variation
by $\sqrt{1/56}<1/4$.  Passing to any function of the history cannot
increase total variation: each event in the output space has a
preimage event in the history space.  Independent randomization also
cannot increase it, by integration over the common seed law.
Consequently \eqref{eq:p-two-point} proves
\begin{equation}\label{eq:p-history-floor}
 R^*(M,h)\ge\frac{3\delta}{8}
 =\frac3{128\sqrt{M+2}}\ge\frac1{64\sqrt M}.
\end{equation}
The last inequality is equivalent to $9M\ge4(M+2)$, valid here.
These IID tables are legitimate all-state witnesses: put
$\theta_u(1)=p$ for every row and initialize the pair with
$\Pi=\operatorname{Bern}(p)^{\otimes2}$.  After two transitions,
$P^2(u,v)=\Pi(v)$ for every initial pair $u$, so the required
certificate $1/M$ holds on every row.

\subsection{Combining the regimes and interpreting precision}

\begin{proof}[Proof of Theorem~\ref{thm:bank-precision}]
For $M<C_0Q^2$, \eqref{eq:p-history-floor} gives
\[
 R^*(M,h)>\frac1{64\sqrt{C_0}\,Q}
 \ge\frac1{1600000Q},
\]
since $\sqrt{C_0}<25000$.
For $M\ge C_0Q^2$, \eqref{eq:p-large-regime} gives this same weaker
lower bound.  Together with \eqref{eq:p-history-floor}, both of
which hold for every decoder, this proves
\[
 R^*(M,h)\ge\frac1{1600000}
              \max\{M^{-1/2},Q^{-1}\}.
\]
The parameter pair used for the statistical floor need not equal the
pair used for the precision floor.  Each is a lower bound for every
decoder's source supremum, so their maximum is a lower bound for the
same minimax risk.

For the upper bound, apply \eqref{eq:a-rounded-risk} with $k=0$ and
$\xi=0$ to exactly the bank in \eqref{eq:p-bank}.  Its sequential
clipping and ratio gate are permitted arbitrary decoders, and
\[
 R^*(M,h)\le\min\left\{1,
 \frac{131072}{\sqrt M}+\frac{536088}{Q}\right\}
 \le\min\left\{1,667160\max\{M^{-1/2},Q^{-1}\}\right\}.
\]
This proves the complete finite-domain statement.
\end{proof}

For any fixed $K>0$, attaining $R^*(M,h)\le K/\sqrt M$ requires
\[
 Q\ge\frac{\sqrt M}{1600000K},\qquad
 h\ge\tfrac12\log_2 M-\log_2(1600000K).
\]
Conversely, $h\ge\lceil\tfrac12\log_2 M\rceil$ suffices by
\eqref{eq:a-rounded-risk}; a separate finite output stage with error
at most $2^{-b}$ adds just that error.  The necessity statement is
about fractional precision of this deterministic nearest-rounded
$J/C/D$ bank.  It is not a lower bound for arbitrary encoders using
$12(h+1)$ bits, arbitrary positional summaries, neural architectures,
or total arithmetic precision.  It gives neither a noise lower bound
nor a $k/M$ lower bound for this changed bank.  Accumulation precision,
addresses, raw input storage, score precision, parameters, masks, and
other program resources are not lower bounded here.  Additional hidden
history or counts would change the premise.  The deliberately loose
crossover $C_0$ is a proof device, not a practical sample-size threshold.

\section{A finite positional attention realization}\label{sec:implementation-proof}

This section realizes the approximate-statistic interface by an explicitly
specified finite causal attention/Boolean--ReLU program. Its purpose is to
verify the observation cut and its costs. A sequential counter on the same
input is a stronger computational baseline, described in
Section~\ref{sec:impl-comparisons}. Uniform aggregation of indicators and
positional copying are established operations
\citep{WeissGoldbergYahav2021,EdelmanEtAl2024}. Moreover, the earlier
finite-bank construction \citep{BCJ001} already contains a paid uniform
positional bank, a real-error half-gate analysis, and a finite sequential
evaluator. The construction here specifies singleton-copy heads, direct
$J/C/D$ rails, nearest rounding, a canonical post-rounding replacement cut,
and an explicit Boolean--ReLU suffix with a separately stated hybrid shape.
It does not establish a new statistical precision phenomenon or a learned
mechanism.

\subsection{The finite model and its guarantee}\label{sec:impl-statement}

For this section only, let $M\ge1$ and $h,b\ge0$ be integers. The parameter
$b$ denotes \emph{output} precision; $r$ remains the raw-window width used
elsewhere. Set
\begin{equation}\label{eq:impl-scales}
 n=M+3,\qquad q=M+2,\qquad Q=2^h,\qquad
 \ell_a=\left\lceil\log_2(n+1)\right\rceil,\qquad
 p=\ell_a+b+h+32.
\end{equation}
The address width $\ell_a$ is distinct from the logarithmic competitive
factor $\ell$ in Theorem~\ref{thm:finite-window}. The input consists of the
$M+2$ binary symbols $Z_0,\ldots,Z_{M+1}$ and one synthetic zero token at
position $q$. No one-hot representation, endpoint label, local pattern
feature, or count is supplied without construction.

For $u=(a,c)\in\{0,1\}^2$, define on an arbitrary binary word
\begin{align}
 J'_u&=\sum_{i=0}^{M-1}\ind\{(Z_i,Z_{i+1})=u\},\notag\\
 C'_u&=\sum_{i=0}^{M-1}\ind\{(Z_i,Z_{i+1})=u,\ Z_{i+2}=1-a\},\label{eq:impl-counts}\\
 D'_u&=\sum_{i=0}^{\lfloor M/2\rfloor-1}
                \ind\{(Z_i,Z_{i+1})=u\}.\notag
\end{align}
An empty sum is zero. A \emph{finite attention head} in the present model
has explicit public structural masks and exact bounded-word evaluation of
its allowed-key sum, followed by a specified rounding operation. The
integer work of that evaluator is charged separately from the number of
attention layers. Pointwise maps consist of affine maps and ReLUs, with
exact finitely represented coefficients and explicit live-wire holds. This
is a hybrid computational model, not unit-cost real arithmetic or a
finite-precision hardware softmax claim.

\begin{theorem}[Finite realization of the positional bank]\label{thm:finite-realization}
For each $M,h,b$ above there is a fixed causal finite program, depending
only on these public parameters, with the following properties.
\begin{enumerate}
\item It has two attention layers and three heads in total, with at most
 two heads in one layer, query/key dimension one, maximum value dimension
 twelve, and score cap zero. Shared pointwise Boolean--ReLU maps have
 depth $O(p^2)$ and width $O(p)$ per token. The integer attention-evaluator
 work is additional and is charged below.
\item On every binary input word, it acquires the exact pair
 $(Z_M,Z_{M+1})$ and produces twelve master numerators
 $B_{u,A}\in\{0,\ldots,Q\}$, $A\in\{J,C,D\}$, satisfying
 \begin{equation}\label{eq:impl-bank-error}
 \left|\frac{B_{u,J}}Q-\frac{J'_u}M\right|,
 \quad\left|\frac{B_{u,C}}Q-\frac{C'_u}M\right|,
 \quad\left|\frac{B_{u,D}}Q-\frac{D'_u}M\right|
 \ \le\ 2^{-h-1}.
 \end{equation}
 The inequalities are simultaneous and deterministic.
\item At the specified cut, the only data-dependent downstream inputs
 are those twelve numerators and the exact endpoint pair. Replace all
 twelve masters jointly by arbitrary canonical integers
 $B^*_{u,A}\in\{0,\ldots,Q\}$ satisfying
 \begin{equation}\label{eq:impl-port}
 \max_{u,A}\frac{|B^*_{u,A}-B_{u,A}|}{Q}\le\delta.
 \end{equation}
 The resulting normalized coordinates have error at most
 $\eta=2^{-h-1}+\delta$ from \eqref{eq:impl-counts} divided by $M$.
 Replacements may be correlated and may depend on the observed word.
 Endpoint bits, metadata, comparators, and subsequent arithmetic remain
 exact; this is the entire numeric fault port.
\item The readout selects the actual endpoint row, clips its $C,D$
 coordinates to $[0,J]$, and applies the gate $J>0$, $512D\ge J$.
 Its output is a binary probability distribution with denominator $2^b$,
 within $2^{-b}$ total variation of the corresponding exact rational gate
 output, including empty rows, gate equality, $h=0$, and $b=0$.
\item In the experiment of Theorem~\ref{thm:approximate-gate}, with
 $M\ge64$ and supplied certificate $1/M$, its actual-terminal risk obeys
 \begin{equation}\label{eq:impl-risk}
 R\le\min\left\{1,\frac{131072}{\sqrt M}
       +6433056\frac{k}{M}
       +1072176\bigl(2^{-h-1}+\delta\bigr)+2^{-b}\right\}.
 \end{equation}
 The program is not supplied $k$ or $\delta$.
\end{enumerate}
A complete sparse-list/sparse-circuit hybrid realization has conservative
online bitwork $O(np^4)$ and total online peak storage
$O(p^4+n\ell_a+np^2)$ bits. A complete dense-mask/dense-circuit alternative
has online bitwork $O(n^2p^2+np^6)$ and total peak storage
$O(p^5+n^2+n\ell_a+np^2)$ bits. These are alternative full
representations, not interchangeable entries in one cost calculation.
Section~\ref{sec:impl-ledgers} gives construction, input, addressing,
mask, hold, cleanup, and output costs and a separately charged all-ReLU
unrolling.
\end{theorem}

The endpoint in this theorem is computed from the input, not appended as
an additional observation. The fault restriction in
\eqref{eq:impl-port} is substantive: off-grid analog errors, corrupted
endpoint bits, accumulator errors, and comparator faults are not covered.
The construction and the following proof do not assume that the rounded
or replaced bank is the count table of any binary word.

\subsection{Causal masks and the twelve direct rails}\label{sec:impl-prefix}

Write $\rho(x)=\max\{x,0\}$. Initially
\[
 \operatorname{raw}_t=Z_t\quad(0\le t\le M+1),\qquad
 \operatorname{raw}_q=0.
\]
Public position information provides the query-type flag
$U_t=\ind\{t=q\}$ and the half flag
\begin{equation}\label{eq:impl-half-flag}
 F_t=\ind\{2\le t<2+\lfloor M/2\rfloor\}.
\end{equation}
These flags are generated by public integer comparisons. Positions and
flags cost $O(n\ell_a)$ stored bits and their generation is charged below.

\paragraph{First layer: two singleton copies.}
Both heads have $Q_t=K_t=0$ and the one-dimensional value
$\operatorname{raw}_t$. Their permitted-key sets are
\begin{equation}\label{eq:impl-copy-masks}
 \mathcal M^{(1)}_t=\begin{cases}\{t-1\},&t\ge1,\\\{t\},&t=0,
                       \end{cases}
 \qquad
 \mathcal M^{(2)}_t=\begin{cases}\{t-2\},&t\ge2,\\\{t\},&t<2.
                       \end{cases}
\end{equation}
Each row is nonempty and causal. Its only normalized attention weight is
exactly one, so the returned bits, denoted $l_{1,t}$ and $l_{2,t}$,
are exact copies. Raw bits and both copied bits are retained through paid
identity holds. In particular,
\begin{equation}\label{eq:impl-copy-identities}
 (l_{2,i+2},l_{1,i+2},\operatorname{raw}_{i+2})
       =(Z_i,Z_{i+1},Z_{i+2})\quad(0\le i<M),\qquad
 (l_{2,q},l_{1,q})=(Z_M,Z_{M+1}).
\end{equation}
The fallback copies at positions zero and one never enter the query's
statistic. These masks are exact structural exclusions. They are not
approximations obtained by assigning large scores in a dense causal row.

\paragraph{Local pointwise features.}
For each $u=(a,c)$ in lexicographic order, form
\begin{align}
 R_{t,u}&=\bigl(1-(l_{2,t}\mathbin{\mathrm{XOR}}a)\bigr)
              \mathbin{\mathrm{AND}}
              \bigl(1-(l_{1,t}\mathbin{\mathrm{XOR}}c)\bigr),\notag\\
 V_{t,u,J}&=R_{t,u},\qquad
 V_{t,u,C}=R_{t,u}\mathbin{\mathrm{AND}}
                    (\operatorname{raw}_t\mathbin{\mathrm{XOR}}a),
 \qquad
 V_{t,u,D}=R_{t,u}\mathbin{\mathrm{AND}}F_t.
 \label{eq:impl-features}
\end{align}
These are twelve binary rails, not an eight-successor histogram with an
additional reconstructed row sum. On the legal key $t=i+2$, substitution
in \eqref{eq:impl-copy-identities} gives respectively
\begin{equation}\label{eq:impl-feature-identities}
 \ind\{(Z_i,Z_{i+1})=u\},\qquad
 \ind\{(Z_i,Z_{i+1})=u,\ Z_{i+2}=1-a\},\qquad
 \ind\{(Z_i,Z_{i+1})=u,\ i<\lfloor M/2\rfloor\}.
\end{equation}
Thus the $C$ rail uses precisely the designated successor. Each feature
uses only its own token's local copies and the public half flag. In
particular, no endpoint-conditioned query is broadcast to historical
tokens.

\paragraph{Second layer: one uniform head.}
This head again has one-dimensional $Q=K=0$, and its value is the
twelve-dimensional vector in \eqref{eq:impl-features}. Its mask is
\begin{equation}\label{eq:impl-uniform-mask}
 \mathcal M^{(3)}_q=\{2,3,\ldots,M+1\},\qquad
 \mathcal M^{(3)}_t=\{t\}\quad(t\ne q).
\end{equation}
The query has exactly $M$ allowed keys. The key $M+1$ includes the last
legal departure triple, whereas the synthetic query $q$ is excluded.
The terminal pair is not an additional departure. Every other row uses a
causal singleton self mask; all those rows are evaluated and their
outputs are subsequently cleared. This choice is part of the construction
and is not a dense causal-prefix restriction.

All allowed scores equal zero, so $\exp(0)=1$ and the ideal normalized
head is a uniform average. There is no approximation of a logarithm,
exponential, anchor inversion, or content match. For each row with $L$
allowed keys, sum each binary rail into an \emph{exact} integer
$S\in\{0,\ldots,L\}$ and return its nearest grid numerator
\begin{equation}\label{eq:impl-round}
 \mathcal R_{h,L}(S)
 =\left\lfloor\frac{2QS+L}{2L}\right\rfloor.
\end{equation}
The denominator $L$ is public, obtained from the enumerated mask, and is
never zero. The floor in this formula will be implemented by the finite
division circuit below; it is not a primitive real-number operation.
At $q$, summing \eqref{eq:impl-feature-identities} proves the all-word
identity
\begin{equation}\label{eq:impl-exact-sums}
 (S_{q,u,J},S_{q,u,C},S_{q,u,D})=(J'_u,C'_u,D'_u).
\end{equation}
Equation~\eqref{eq:impl-round} rounds $QS/L$ to its nearest integer, with
ties upward. It maps $0$ to $0$ and $L$ to $Q$, and
\[
 \left|\frac{\mathcal R_{h,L}(S)}Q-\frac SL\right|
 \le\frac1{2Q}.
\]
Together with \eqref{eq:impl-exact-sums}, this proves
\eqref{eq:impl-bank-error} for every binary word, including $M=1$ and
$h=0$. Counts are accumulated before this single rounding. Repeated
$h$-grid rounding during accumulation would be a different algorithm.

\subsection{Exact Boolean and ReLU arithmetic}\label{sec:impl-arithmetic}

All words use little-endian bit rails, with enough leading zeros for the
stated operation. For bits $x,y$, the identities
\begin{align}
 x\mathbin{\mathrm{AND}}y&=\rho(x+y-1),&
 \mathrm{NOT}(x)&=\rho(1-x),\notag\\
 x\mathbin{\mathrm{XOR}}y&=\rho(x-y)+\rho(y-x),&
 x\mathbin{\mathrm{OR}}y&=x+y-\rho(x+y-1)
 \label{eq:impl-boolean}
\end{align}
are exact. Each expression is a constant-size affine/ReLU circuit. A
coordinatewise multiplexer is
\begin{equation}\label{eq:impl-mux}
 \operatorname{mux}(s,x,y)
 = (s\mathbin{\mathrm{AND}}x)
      \mathbin{\mathrm{OR}}
   ((\mathrm{NOT}(s))\mathbin{\mathrm{AND}}y).
\end{equation}
It equals $x$ when $s=1$ and $y$ when $s=0$. Thus no data-dependent
selection below is free dynamic addressing.

\paragraph{Addition, subtraction, and comparison.}
For input bits $x_i,y_i$ and carry $c_i$, a full adder has
\[
 v_i=x_i\mathbin{\mathrm{XOR}}y_i,\qquad
 s_i=v_i\mathbin{\mathrm{XOR}}c_i,\qquad
 c_{i+1}=(x_i\mathbin{\mathrm{AND}}y_i)
             \mathbin{\mathrm{OR}}
            (v_i\mathbin{\mathrm{AND}}c_i).
\]
Induction on $i$ gives ordinary binary addition. A ripple chain on
$w$-bit words uses $O(w)$ depth and active gates, with $O(w)$ live rails.
Retain the final carry separately. Subtract $y$ from $x$ by adding
$x+\mathrm{NOT}(y)+1$. Since this integer is $x-y+2^w$, its final carry
is one exactly when $x\ge y$, and its low $w$ bits are $x-y$ modulo
$2^w$. Hence the same chain implements unsigned comparison and
subtraction. Public shifts are explicit fixed wires with zero padding;
their storage and holds are charged.

\begin{lemma}[Finite restoring division]\label{lem:impl-division}
For a positive $w$-bit denominator and a $w$-bit dividend, exact quotient
and remainder can be computed with $w$ fixed public iterations, $O(w^2)$
active Boolean/ReLU gates and depth, and $O(w)$ live rails. A layered
representation that holds all live rails at every step has at most
$O(w^3)$ gates and identity holds.
\end{lemma}
\begin{proof}
Initialize remainder $R=0$. Scan the dividend's bits from most significant
to least. Given the next bit $x_i$, first form $\widetilde R=2R+x_i$.
Subtract the denominator $d$ using the preceding circuit. Its final carry
$c$ says whether $\widetilde R\ge d$. Emit $c$ as the next quotient bit,
and replace the remainder by
$\operatorname{mux}(c,\widetilde R-d,\widetilde R)$.

After scanning any prefix $P$ of the dividend, the invariant is
$P=Q_Pd+R$ with $0\le R<d$, where $Q_P$ is the emitted quotient prefix.
For the next bit, $2R+x_i<2d$, so one subtraction suffices, proving the
invariant inductively. Also $R\le P$, and therefore
$\widetilde R\le2P+x_i$, the next scanned dividend prefix. That prefix
fits in the original $w$ bits. Thus the shifted remainder fits without an
uncharged extra magnitude bit. The subtraction's carry is one explicit
additional rail. At the final iteration the invariant is the desired
Euclidean division identity. Each of $w$ iterations costs $O(w)$ active
gates and depth. Holding $O(w)$ live rails through $O(w^2)$ layers gives
the stated conservative $O(w^3)$ bound.
\end{proof}

The same sequence is a total Boolean circuit even on denominator zero;
its arbitrary zero-denominator result will never be used. The readout
below instead selects a strictly positive denominator before invoking
this circuit.

\paragraph{Word bounds for the attention evaluator.}
A feature sum needs $\lceil\log_2(M+1)\rceil$ bits; $h$ bits alone do not
suffice. The rounding dividend and divisor satisfy
\begin{equation}\label{eq:impl-prefix-wordbounds}
 2QS+L\le 2^{h+1}M+M<2^{h+\ell_a+2},\qquad
 2L\le2M<2^{\ell_a+1}.
\end{equation}
The same bounds cover singleton rows because $M\ge1$. Use
$w_{\rm att}=\ell_a+h+16$ bits, together with explicitly retained carries.
The accumulation, the $(h+1)$-place shift, addition of $L$, and restoring
division all fit. In particular, the accumulation is exact, and the
rounding numerator is computed without overflow. This implements
\eqref{eq:impl-round} entirely by the preceding finite gadgets.

\subsection{Packed and scalar versions of the bank}\label{sec:impl-extraction}

The bank has two possible physical representations, with the same declared
observation. In a packed version, each master is stored as its $h+1$
numerator bits. Reading these existing bits and padding them with zeros
is paid fixed wiring. The high bit is necessary to represent $B=Q$, that
is, the normalized value one.

In a scalar version, the head returns one exact grid scalar
$H=B/Q$ per rail. The head's conversion from numerator bits to this
scalar is the paid dyadic sum
\[
 H=\sum_{j=0}^{h}2^{j-h}B_j.
\]
It uses $O(h+1)$ terms with $O(h+1)$-bit coefficients. A scalar is
\emph{not} treated as exposing its bits at no cost.

\begin{lemma}[Paid extraction from a canonical scalar grid]
\label{lem:impl-extraction}
For $H\in\{0,1/Q,\ldots,1\}$, the numerator bits of $QH$ are extracted
exactly by a ReLU circuit of depth and active-gate count $O(h+1)$, with
$O((h+1)^2)$ conservative holds and $O(h+1)$-bit coefficients and scalar
words. No floor on arbitrary real inputs is used.
\end{lemma}
\begin{proof}
Initialize $A=QH$ and, in the fixed order $j=h,h-1,\ldots,0$, set
\begin{equation}\label{eq:impl-extraction}
 B_j=\rho\bigl(A-(2^j-1)\bigr)-\rho(A-2^j),\qquad
 A\leftarrow A-2^jB_j.
\end{equation}
Initially $A$ is an integer in $[0,Q]$. At the beginning of step $j$ it
is an integer in $[0,2^{j+1}-1]$. If $A<2^j$, both ReLUs in
\eqref{eq:impl-extraction} vanish. If $A\ge2^j$, their difference is
exactly one. The update leaves an integer in $[0,2^j-1]$. Induction gives
all the bits and final remainder zero. At $H=1$, the top bit $B_h$ is
one and all remaining bits are zero. At $h=0$, the single step extracts
either zero or one. Each step has constant active cost; retaining all
emitted bits through later steps costs $O((h+1)^2)$ identity holds.
Thresholds, scaling coefficients, and exact scalar values have
$O(h+1)$-bit descriptions.
\end{proof}

For the scalar realization the replacement supplies a canonical scalar
$H^*=B^*/Q$, and extraction occurs \emph{after} replacement. Consequently
it reveals only the replaced master, not its earlier value. On off-grid
inputs the recurrence need not produce bits; such inputs are outside the
fault model. Twelve parallel extractions, their scaling, their scalar
head decodes, and their holds fit the $p$-based resource bounds below.
Both representations therefore reach the same padded numerator rails
by paid operations.

\subsection{Endpoint selection and a total finite readout}
\label{sec:impl-readout}

Let $(e_0,e_1)=(Z_M,Z_{M+1})$ be the two copied endpoint bits. For each
$u=(a,c)$ form the selector
\[
 s_u=\bigl(1-(e_0\mathbin{\mathrm{XOR}}a)\bigr)
          \mathbin{\mathrm{AND}}
       \bigl(1-(e_1\mathbin{\mathrm{XOR}}c)\bigr).
\]
Exactly one selector equals one. For each of the three coordinate types,
start a zero word and traverse the four rows in fixed order, replacing
it by $\operatorname{mux}(s_u,B^*_{u,A},\text{current word})$.
The resulting words are the selected canonical numerators. This is a
constant number of word muxes; it does not use an uncharged lookup indexed
by the endpoint.

Denote the selected words by $J,C_0,D_0$, and set
\begin{equation}\label{eq:impl-clip-gate}
 C=\min\{C_0,J\},\qquad D=\min\{D_0,J\},\qquad
 G=\ind\{J\ge1\}\mathbin{\mathrm{AND}}\ind\{512D\ge J\}.
\end{equation}
Each minimum is implemented by comparison and a mux. Lower clipping is
automatic from unsigned canonical inputs. Multiply $D$ by $512$ with a
nine-place shift and compare exactly. Gate equality therefore accepts,
whereas $J=0$ rejects, irrespective of the second comparison.

Choose the safe denominator and quotient by
\begin{equation}\label{eq:impl-output}
 J_{\rm safe}=\operatorname{mux}(G,J,1),\qquad
 v=\left\lfloor\frac{2^bC}{J_{\rm safe}}\right\rfloor,\qquad
 m=\operatorname{mux}(G,v,0).
\end{equation}
The divisor is always positive. The shift and restoring division implement
the displayed floor. Even on a rejected row the pre-mux quotient is at
most $2^{h+b}$. On acceptance, $0\le C\le J$ gives $0\le m\le2^b$.
On rejection $m=0$. Compute the complement numerator $2^b-m$ by exact
subtraction and route the designated probability using the endpoint bit:
\begin{equation}\label{eq:impl-symbol-route}
 (P_0,P_1)=
 \begin{cases}
 (2^b-m,m),&e_0=0,\\
 (m,2^b-m),&e_0=1.
 \end{cases}
\end{equation}
Again this is a pair of bitwise muxes. The actual output is
$(P_0/2^b,P_1/2^b)$, always in the binary simplex. A dyadic decoder
$\sum_{i=0}^b 2^{i-b}(P_a)_i$ uses paid coefficients of $O(b+1)$ bits.

On acceptance, flooring gives
$0\le C/J-m/2^b<2^{-b}$; on rejection the rational and finite outputs
both put zero on the designated symbol. Binary total variation is the
absolute difference of that probability. This proves the output error
bound, including $b=0$, without a gate-margin assumption.

All readout arithmetic can use $w_{\rm out}=h+b+16$ bits. Indeed,
$512D\le2^{h+9}$, $2^bC\le2^{h+b}$, the pre-mux quotient is no larger
than this dividend, and the final probability numerators are at most
$2^b$. Constants, intermediate carries, and padded words are bounded by
$2p$ rails. Together with \eqref{eq:impl-prefix-wordbounds}, these bounds
prove the absence of overflow throughout the construction.

The pointwise feature and suffix maps are shared and executed at
\emph{all} $n$ tokens. The public flag $U_t$ clears packed bank and endpoint
inputs at nonquery tokens using muxes. For a scalar bank rail
$H_t\in[0,1]$, the paid gate $\rho(H_t+U_t-1)$ retains $H_t$ at the
query and is zero elsewhere. Thus nonquery banks are cleared before the
cut without extracting the query's bits early. Scalar extraction after
replacement is executed at all tokens, with zero inputs away from the
query. The shared suffix still computes on those zero inputs; its work
is not charged only at $q$. Only the query's output is retained.

\subsection{Isolation of the replacement cut and statistical composition}
\label{sec:impl-cut}

The master cut is after the twelve query coordinates have been rounded
and before endpoint-conditioned selection, comparison, division, or any
readout fanout. Every downstream bank copy originates from the replaced
master. In the scalar realization it also precedes the extraction in
Lemma~\ref{lem:impl-extraction}. The endpoint bits and public
$h,b$, type flags, and layout remain exact. Nonquery banks and endpoint
rails are cleared by the public query flag.

Before the cut is made available to the suffix, raw-input arrays, lag-copy
arrays, local-feature arrays, attention accumulators, rounding dividends
and divisors, intermediate quotients and remainders, and any clean-master
duplicates are released or overwritten. The suffix has no later attention
layer and no raw-memory read. Registers retaining $M$ or layout contain
only public, data-independent values; they do not encode any discarded
count. An equivalent interface signature is simply
\[
 \bigl((B^*_{u,A})_{u,A},\ (e_0,e_1),\ h,b\bigr)
       \longmapsto (P_0,P_1).
\]
There is no original-bank handle in this signature. All initialization,
clearing, and identity holds are included in the resource accounts.

This is a logical access restriction on the program, not a claim of
secure physical erasure against a hostile host retaining its own copies.
The external fault injector is not counted as free neural hardware: its
construction and runtime are outside the nominal program. The theorem
charges the program on either side of the stated replacement port, not
an implementation of the adversary.

\begin{proof}[Completion of the proof of Theorem~\ref{thm:finite-realization}]
The copied-bit identity \eqref{eq:impl-copy-identities}, the feature
identity \eqref{eq:impl-feature-identities}, and the exact mask
\eqref{eq:impl-uniform-mask} prove the twelve all-word sums
\eqref{eq:impl-exact-sums}. Exact finite division implements their nearest
rounding, proving \eqref{eq:impl-bank-error}. The replacement triangle
inequality yields
\[
 \left|\frac{B^*_{u,A}}Q-\frac{A'_u}M\right|
 \le\frac1{2Q}+\delta=\eta
 \qquad(A\in\{J,C,D\}).
\]
For analysis only, multiply the normalized replacement coordinates by
$M$, obtaining approximate counts with simultaneous error at most
$e=M\eta$. No such multiplication by $M$ is needed or supplied to the
actual readout. The canonical $J$ already lies in $[0,Q]$, and
\eqref{eq:impl-clip-gate}, after the common positive scaling $M/Q$,
is exactly the sequential clipping and gate of
Theorem~\ref{thm:approximate-gate}. Both the ratio and the comparison are
unchanged by this scaling. Inconsistent replacement coordinates are
permitted because that theorem does not require count realizability.

The protected endpoint equals the actual terminal pair and has been
acquired by the first layer. Equations~\eqref{eq:impl-output}--
\eqref{eq:impl-symbol-route} give output error at most $\xi=2^{-b}$.
Applying Theorem~\ref{thm:approximate-gate} with $e=M\eta$ and that $\xi$
proves \eqref{eq:impl-risk}. This is the only stochastic step. The
all-word interface identities need neither a source law nor an attack
independence assumption. The statistical composition retains the source,
strict-edit, and observation-measurability assumptions of the approximate
gate theorem. The finite circuits above give the claimed shape; the next
subsection verifies their full cost bounds.
\end{proof}

In particular, if $h,b\ge\lceil\tfrac12\log_2M\rceil$ and
$\delta=O(M^{-1/2})$, \eqref{eq:impl-risk} is
$O(M^{-1/2}+k/M)$. Theorem~\ref{thm:bank-precision} calibrates the clean,
fault-free \emph{fixed bank} obtained at this cut. It does not make the
present word widths or total program precision necessary. The phrase
``half-log precision'' can refer only to the fractional precision $h$ of
these normalized coordinates, not to all arithmetic, weights, addresses,
input storage, or total memory.

\subsection{Complete resource accounts}\label{sec:impl-ledgers}

All following bounds are constructive asymptotic upper bounds from the
gate recipes and fixed public schedules, with absolute constants. They
are not timing measurements or assertions that a conventional softmax
kernel executes the integer evaluator at unit cost. We count a bit
operation at the elementary Boolean/arithmetic level and charge the
selected program representation, addresses, and all holds. Only one
representation is selected at a time.

\paragraph{Hybrid shape and coefficient descriptions.}
There are two attention layers: two one-bit copying heads followed by one
twelve-value uniform head. All scores are zero; query/key dimensions are
one. A constant number of $O(p)$-word selections, comparisons, and
restoring divisions gives shared pointwise depth $O(p^2)$ and width
$O(p)$. The twelve scalar extractions, when present, add only $O(h+1)$
depth and fit that width. Attention's exact integer evaluator is charged
as work below and is not hidden in this small hybrid depth.

Individual coefficients use $O(p)$ bits. Boolean coefficients are bounded
constants, scalar-extraction coefficients use $O(h+1)$ bits, and output
dyadics use $O(b+1)$ bits. The shared pointwise circuit has $O(p^2)$ active
gates. Explicitly holding $O(p)$ live rails through all layers gives
$O(p^3)$ sparse gate/hold records. Their addresses and coefficients cost
at most $O(p)$ bits each, hence $O(p^4)$ stored sparse parameter bits.
The dense alternative has $O(p^2)$ layers, $O(p)$ width, and therefore
$O(p^4)$ coefficient slots including zeros and biases. Storing all
coefficients costs $O(p^5)$ bits. Constantly many head-projection matrices
fit these bounds. Finite evaluator code and bounded-word templates cost
at most $O(p^4)$ bits and are included, not left as an external oracle.

\paragraph{Masks, public positions, and input.}
Binary input acquisition costs $M+2$ bits. The synthetic zero is public.
For positions, query type, half flags, row boundaries, and masks, the two
complete choices are:
\begin{enumerate}
\item \emph{Active-key lists.} The first two heads each contain $n$
 stored keys. The third contains $n-1+M$ keys. Thus the exact number of
 stored key occurrences is $3n-1+M$. Each key and row boundary uses
 $O(\ell_a)$ bits, giving $O(n\ell_a)$ mask/list storage, including
 positions and flags. Enumerating indices, comparisons, routes, and row
 boundaries costs at most $O(n\ell_a^2)$ bitwork. Each stored route is
 read and charged during execution. This is explicitly paid sparse
 structural masking, not free addressing.
\item \emph{Literal dense masks.} Three $n$-by-$n$ masks cost exactly
 $3n^2$ mask bits, in addition to $O(n\ell_a)$ metadata. Generating and
 writing them costs at most $O(n^2\ell_a)$ bitwork. A dense online scan
 costs at most $O(n^2p)$ even though many permitted-key sets are
 singletons. Converting active lists to this representation pays these
 dense generation and storage costs.
\end{enumerate}

\paragraph{Finite attention arithmetic.}
There are $O(n)$ allowed key occurrences in the active-list version.
Each copied value is a bit. For the uniform head, each of twelve rails
per key uses an $O(p)$-bit addition, and every token performs twelve
$O(p^2)$ rounding divisions. Thus the fixed-mask finite attention
evaluation costs $O(np^2)$ elementary bitwork, including copies,
projections, and clearing. No logarithm, exponential, or source-specific
cache is generated. Dense-mask scanning adds the cost just stated.
Controllers and word addresses occupy $O(p)$ bits; the larger complete
runtime bounds below also allow a full $O(p)$ overhead per addressed
record lookup.

\paragraph{Two complete online implementations.}
Sparse execution of all pointwise maps, including every live hold, costs
$O(np^3)$ elementary bitwork. Combining this with active-list attention,
then charging up to $O(p)$ per sparse record access, gives the conservative
unconditional sparse/sparse online bound
\begin{equation}\label{eq:impl-sparse-work}
 T_{\rm sparse}=O(np^4).
\end{equation}
Its stored parameters and code are $O(p^4)$ bits, its masks and public
metadata are $O(n\ell_a)$ bits, and its observed input costs $M+2$ bits.
Dense coefficient counts are not substituted into this runtime.

For the dense-mask/dense-circuit alternative, coefficient-bit operations
are charged even for zero terms. Dense QKV projections, all-token shared
maps, mask scans, and word operations give the conservative bound
\begin{equation}\label{eq:impl-dense-work}
 T_{\rm dense}=O(n^2p^2+np^6).
\end{equation}
Its parameters and code cost $O(p^5)$ bits, with masks and metadata
$O(n^2+n\ell_a)$ bits. The quadratic term is a safe allowance for this
dense implementation, not a lower bound or an assertion that the sparse
implementation performs quadratic work.

\paragraph{Activations, cleanup, and total online peak.}
With bit rails and packed words, the working activation arrays require
$O(np)$ bits. A conservative scalar-channel implementation allocates
$O(p)$ physical bits to each of $O(p)$ live channels per token, hence
$O(np^2)$ bits. These allowances include Q/K/V arrays, feature residuals,
lag holds, twelve pre-clearing bank words per token, endpoint holds,
controller, and arithmetic scratch. Input costs $O(n)$ additional bits;
output costs $O(b+1)$ packed bits, or $O((b+1)p)$ in the conservative
scalar allocation. The explicit +1 covers $b=0$ and the numerator $2^b$.
Initializing, overwriting, and releasing these arrays is charged in the
online bounds. Only the latest live layer is retained: a training tape or
all-layer retention would be an additional cost.

Total online peak is the sum of simultaneously resident program and
parameters, masks and metadata, observed input, QKV and residual arrays,
banks, controller and scratch, and output. With conservative scalar
allocation it is bounded by
\begin{align}
 S_{\rm sparse}&=O(p^4+n\ell_a+np^2),
                  \label{eq:impl-sparse-memory}\\
 S_{\rm dense}&=O(p^5+n^2+n\ell_a+np^2).
                  \label{eq:impl-dense-memory}
\end{align}
The $12(h+1)$ bank bits and two endpoint bits are only the cut payload,
not these total-memory quantities. At $h$ of order $\tfrac12\log_2M$,
exact accumulators, addresses, and public horizon words still require
order $\log M$ bits, and the resident input and paid mask storage remain.

\paragraph{Offline construction and peak storage.}
Generate the ripple circuits, fixed division schedules, muxes, scalar
extractions, and all identity holds in $O(p^5)$ conservative bitwork with
$O(p^4)$ sparse scratch/record storage, and write the $O(p^4)$ sparse
program bits. Sorting, if used, involves $O(p^3)$ records of $O(p)$ bits
and $O(\log p)$ comparison overhead, which fits the same work allowance.
The dense representation additionally writes $O(p^5)$ coefficient bits,
including zero slots, within $O(p^5)$ bitwork. Add
$O(n\ell_a^2)$ work for sparse lists or $O(n^2\ell_a)$ work for dense
masks and their public metadata. Accordingly, safe full construction
work bounds are
\begin{equation}\label{eq:impl-offline-work}
 O(p^5+n\ell_a^2)\quad\text{(sparse)},\qquad
 O(p^5+n^2\ell_a)\quad\text{(dense)}.
\end{equation}
Offline peak includes partially generated stored program and masks plus
live generation scratch. It is bounded by the corresponding total in
\eqref{eq:impl-sparse-memory} or \eqref{eq:impl-dense-memory}; retaining
both construction and execution workspaces adds their sum, of the same
order. There is no hidden training phase, source-specific compilation,
or external arithmetic oracle.

\paragraph{A separate all-Boolean/ReLU unrolling.}
To remove the finite-attention primitive, select one complete
representation and unroll its fixed public bit-operation schedule. Let
$T$ be that representation's conservative online bitwork and $S$ its total
live bit storage, including input, program/constants, masks, and scratch.
Every conditional data operation is already a Boolean mux; loop lengths
and permitted routes are public. Replace each elementary bit operation by
\eqref{eq:impl-boolean}--\eqref{eq:impl-mux}, and at every instruction
explicitly hold every other live rail. Induction on instructions proves
that the resulting network state equals the original finite evaluator's
state. The final dyadic decode is also retained and charged.

This produces a \emph{PURE} Boolean/ReLU network, with no finite-attention
arithmetic primitive, having the valid bounds
\begin{align}
 \text{depth}&=O(T+p),&
 \text{width}&=O(S+p),\notag\\
 \text{sparse gates and holds}&=O((T+p)(S+p)),&&\label{eq:impl-pure}\\
 \text{dense parameter bits}&=O\bigl(p(T+p)(S+p)^2\bigr),&&\notag\\
 \text{dense bitwork}&=O\bigl(p^2(T+p)(S+p)^2\bigr).&&\notag
\end{align}
For example, one may use
$T=O(np^4)$ and $S=O(p^4+n\ell_a+np^2)$ from the sparse hybrid.
The sparse-record description can conservatively be written with
$O(p(T+p)(S+p))$ bits: all relevant indices have $O(p)$ bits, since
$\log n\le p$ and $\log p\le p$.
Writing and evaluating the PURE representation incurs its own displayed
storage and work costs. These large unrolling bounds do not inherit the
hybrid depth $O(p^2)$ or per-token width $O(p)$. The small hybrid shape
therefore must never be presented as the shape of the fully expanded
finite evaluator.

\subsection{Observation comparisons and the classical baseline}
\label{sec:impl-comparisons}

\paragraph{The count-only obstruction has an information scope.}
Theorem~\ref{thm:finite-window}, specialized to $r=2$, gives an optimal
budget-blind competitive factor between
\[
 2^{-30}\left(1+\log\frac{M+2}{2}\right)
 \quad\text{and}\quad
 2^{20}\left(1+\log\frac{M+2}{2}\right)
\]
for exact triple counts plus both endpoint pairs, at supplied $1/M$.
The lower tests the same map across budgets; it is not a logarithmic
penalty at every individual budget. Counts plus only the protected pair,
and any rounded versions of those observations, are weaker and inherit
the obstruction. The present half-position statistics change the
observation, so that comparison supplies neither an architecture-wide
lower bound nor a neural depth separation.

\paragraph{A one-pass counter on equal positional information.}
On the same attacked word, a classical program keeps four exact triples
of counters $(J_u,C_u,D_u)$ and a two-bit sliding window. When symbol
$Z_{i+2}$ arrives, the window identifies $u=(Z_i,Z_{i+1})$: increment
$J_u$, increment $C_u$ if $Z_{i+2}=1-Z_i$, and increment $D_u$ if
$i<\lfloor M/2\rfloor$. Then shift the two-bit window. After the last
symbol it contains the protected endpoint. Thus the program computes
exactly \eqref{eq:impl-counts} in one pass, without knowing $k$.

Its counters, loop index, and sliding window use $O(\log M)$ auxiliary
bits, with $O(M\log M)$ conservative bitwork for the scan. On the exact
counters, the same positive-denominator check, gate, and finite output
division give the risk in \eqref{eq:impl-risk} with $\eta=0$.
If output precision is charged, the division uses
$O((\log M+b+1)^2)$ active bitwork and
$O(\log M+b+1)$ live bits, with additional holds if fully layered.
Alternatively, round the exact counters once at the end using
\eqref{eq:impl-round}. This reproduces the attention bank coordinate for
coordinate, so it is also an equal-bank baseline, not merely an
asymptotically comparable estimator.

The input models must be matched. If the word is resident, charge the
same $M+2$ input bits to the counter and the attention program. If a
read-once input stream is allowed, state that permission explicitly; do
not compare the counter's auxiliary workspace with the attention
program's batch total memory. The counter has no attention mask-storage
cost. Consequently no neural statistical or computational advantage over
equal positional information is asserted. Sparse masks make the proposed
realization cheaper than generic dense attention evaluation, but do not
surpass this sequential baseline.

\paragraph{What changes relative to the earlier finite bank.}
The construction of \citet{BCJ001} already supplies eight full-triple and
four half-row normalized coordinates, public half-position flags,
uniform attention, endpoint acquisition, an isolated finite observation,
and a paid finite sequential evaluator with optional full unrolling. It
also already distinguishes fractional bank precision from arithmetic word
width and allows a final rounded probability. Its approximate-gate,
quantization, noise, and source-edit results establish the central
finite-bank precision phenomenon in that earlier interface.

The present implementation makes a different, explicit choice at several
points. Two singleton heads acquire the lagged bits and endpoint. The
local Boolean circuit emits four direct $J/C/D$ triples. A zero-score
head uses exact accumulation followed by one nearest rounding, rather
than the earlier evaluator's nested floors. The numeric fault is a
canonical replacement \emph{after} that rounding, rather than an additive
analog change before a final floor. Bit circuits, scalar-grid extraction,
restoring division, and dyadic output give the specified two-layer,
three-head hybrid with growing pointwise depth. Before rounding, the
full-triple and direct-row representations are linearly interconvertible;
after separately applied rounding they need not define the same
statistical experiment. Neither port theorem, especially a port-noise
lower bound, transfers merely because both banks have twelve labels.

Uniform indicator aggregation, positional copying, and finite arithmetic
compilation are standard ingredients
\citep{WeissGoldbergYahav2021,MerrillSabharwal2023,EdelmanEtAl2024,
RajaramanEtAl2024}. The construction here uses exactly accumulated
indicators followed by rounding of their aggregate, not rounding each
attention weight $1/M$ onto the $h$-grid. Results about rounding attention
probability vectors therefore concern a different numerical operation.
The explicit singleton structural masks likewise must not be confused
with dense-causal-only attention or finite-temperature approximations to
copying. No conclusion is asserted about training, learnability, minimal
head count, total-precision optimality, architecture-wide lower bounds,
pure-network depth advantages, empirical performance, or priority for
these standard mechanisms.

\clearpage
\begingroup
\raggedright
\bibliographystyle{plainnat}
\bibliography{references}
\endgroup
\end{document}
